%%% A claim-ledger walkthrough of the SPS confidentiality checker.
%%%
%%% Build:  latexmk -pdf main.tex
%%% Fast:   latexmk -pdf testbuild.tex   (one part at a time)

%%% Preamble for a claim-ledger walkthrough document.
%%%
%%% Copy into the document directory and \input it, or paste inline.
%%% Adapt the theorem environments, title and notation package to the project;
%%% leave the geometry and the commented workarounds alone unless you have
%%% re-diagnosed the problem they solve (see references/authoring-notes.md).

\documentclass[11pt,letterpaper]{article}

%% Wide outer margin: the claim ledger lives there and takes ~1/5 of the page.
%% `includemp` makes the margin column part of the body width, so:
%%   8.5in = 0.85 (left) + 5.35 (text) + 0.13 (sep) + 1.72 (ledger) + 0.45 (right)
\usepackage[letterpaper,
            left=0.85in, right=0.45in, top=0.95in, bottom=1.0in,
            marginparwidth=1.72in, marginparsep=0.13in, includemp]{geometry}

\usepackage{amsmath,amssymb,amsthm}
\allowdisplaybreaks
\usepackage{booktabs}
\usepackage{array}
\usepackage{enumitem}
\usepackage{microtype}
\usepackage{xcolor}

%% float: gives `algorithm` the [H] placement specifier, so a sub-algorithm
%% stays inside its step instead of floating away from the ledger panel it is
%% aligned with.
\usepackage{float}
\usepackage{algorithm}
\usepackage{algpseudocode}
%% algorithmic restarts its line counter per algorithm, so hyperref emits
%% duplicate ALG@line.N destinations. Qualify them by algorithm number.
\makeatletter
\newcommand{\theHALG@line}{\thealgorithm.\arabic{ALG@line}}
\makeatother

\usepackage{tikz}
\usetikzlibrary{positioning, arrows.meta, calc, decorations.pathreplacing}
%% \captionof lets diagrams carry numbered captions without being floats --
%% floats drift away from their step and desynchronise from the ledger.
\usepackage{caption}

\usepackage[colorlinks=true,
            linkcolor=blue!50!black,
            citecolor=blue!50!black,
            urlcolor=blue!50!black]{hyperref}
\usepackage[capitalize,nameinlink]{cleveref}

%% Latin Modern + T1: Computer Modern lacks typewriter braces, bold small caps,
%% and the sub-5pt sizes the ledger annotations ask for. lmodern supplies all
%% three; anyfontsize covers any remaining odd size.
\usepackage[T1]{fontenc}
\usepackage{lmodern}
\usepackage{anyfontsize}
%% NOTE: the outcome tokens are small caps and `theorem' bodies are italic, so a
%% token placed inside a theorem asks for small-caps italic, a shape Latin Modern
%% does not cut. Do not try to fix that with \DeclareFontShape here: lmodern
%% declares its families lazily via .fd files, so a preamble declaration fails
%% with "Font family `T1+lmr' unknown". Keep the tokens out of theorem bodies
%% instead -- see the wording of the diagnostic-narrowing theorem.

\usepackage{claimledger}
%\usepackage{<project>notation}     % <- write one per project

%% Typewriter-heavy paragraphs: cmtt spaces do not stretch, so give the line
%% breaker slack instead of letting it emit overfull lines.
\tolerance=1600
\setlength{\emergencystretch}{1.8em}
\hbadness=2500

\theoremstyle{definition}
\newtheorem{definition}{Definition}[section]
\newtheorem{remark}[definition]{Remark}
\newtheorem{obligation}[definition]{Obligation}
\theoremstyle{plain}
\newtheorem{theorem}[definition]{Theorem}
\newtheorem{lemma}[definition]{Lemma}
\newtheorem{proposition}[definition]{Proposition}

\crefname{definition}{Definition}{Definitions}
\crefname{remark}{Remark}{Remarks}
\crefname{obligation}{Obligation}{Obligations}
%% `step` environments \refstepcounter{clgstep}, so a \label inside one binds to
%% clgstep -- teach cleveref to name it, or \cref emits a warning and \ref
%% silently yields a step number where a section number is implied.
\crefname{clgstep}{Step}{Steps}
\Crefname{clgstep}{Step}{Steps}
%% The document calls its top-level divisions "Part N" throughout, but they are
%% \section units, so a bare \cref would render "Section N" and collide with the
%% "Step N" vocabulary reserved for the walkthrough.
\crefname{section}{Part}{Parts}
\Crefname{section}{Part}{Parts}

\usepackage{spsnotation}

\title{\textbf{Proving a program leaks nothing}\\[2pt]
       \large A step-by-step walkthrough of the SPS confidentiality checker,\\
       \large its evidence pipeline, and its LLVM/MLIR test harness}
\author{Claim-ledger lecture notes}
\date{}

\begin{document}
\maketitle

\begin{abstract}
\noindent
This is a walkthrough of what it takes to answer one question mechanically:
\emph{does this compiled artifact reveal a secret to an observer who is not
entitled to it?} The answer is never a single proof. It is a chain of
reductions, each trading one obligation for a different obligation that is
easier to attack, and the interesting content of the system lives in the
trades rather than in any individual step.

The right margin of every step carries a \emph{claim ledger}: the obligations
still open at that point in the argument, colour-coded by what the step just
did to them. A reader who follows only the margin still learns the shape of
the argument. The last step leaves the ledger empty.

Part~\ref{part:harness} then maps every step onto the regression fixtures that
guard it, and Part~\ref{part:correctness} answers the question that motivates
all of it: how would you know the checker itself is correct, and what would a
mechanized proof for an MLIR-based analysis actually look like?

Part~\ref{part:lowering} then confronts the finding that undoes much of the
rest: the artifact a checker reads is not the artifact that runs. Twenty
reproduced compiler transformations force an architecture in which policy stops
travelling downward and observations travel upward instead.

Part~\ref{part:implementation} is the shortest and the only one about something
that runs. Six claims in the parts above were wrong until a tool was executed
against them, which is the argument for building the smallest real thing before
designing any further.
\end{abstract}

\tableofcontents

\section{Orientation}
\label{part:orientation}

\subsection{The statement}

Fix a compiled artifact $P$. Partition its inputs into public inputs $p$ and
secret inputs $s$. Fix an observer and let $\Obs_\Theta$ be the projection of an
execution trace onto what that observer can actually see under target profile
$\Theta$ --- branch directions, memory addresses, operation latencies and
counts, allocation sizes, transfer lengths, values arriving at public sinks, and
the identity of the host or audience a message is delivered to.

Some dependence on the secret is intended: a service that returns nothing
derived from its inputs is useless. So the property is not plain
noninterference but \emph{release-relative} noninterference. Let $\Rel$ be the
releases authorized for that observer. Then:
%
\begin{equation}
\label{eq:ni}
\forall p,\, s_0,\, s_1 :\quad
  \Rel(p,s_0) = \Rel(p,s_1)
  \;\Longrightarrow\;
  \Obs_\Theta\bigl(\Trace(P,p,s_0)\bigr) = \Obs_\Theta\bigl(\Trace(P,p,s_1)\bigr).
\end{equation}
%
Read it as a challenge from an adversary. They choose the public inputs once and
two secrets freely. They only have to respect one constraint: whatever they are
\emph{entitled} to learn must come out the same both times. If they can then
make any observable differ, they have learned something they were not entitled
to, and $P$ leaks.

\begin{remark}
Equation~\eqref{eq:ni} quantifies over \emph{pairs} of executions. That single
fact drives almost every design decision in the rest of this document, and
Part~\ref{part:background} is devoted to its consequences.

Two warnings about the form above, both discharged there. It is a statement of
\emph{intent}, not an encoding: installing its antecedent as an initial
constraint on the pair is a known unsoundness, and \cref{sec:ledger} gives the
prefix-causal construction that replaces it. And $P$ here is informal shorthand;
from \cref{step:statement} onward the object under analysis is the frozen
artifact $\artifact$ together with its pipeline $\pipeline$, because a verdict is
only ever a statement about specific bytes.
\end{remark}

\subsection{Why functional correctness is not enough}

A verified compiler such as CompCert proves that the target program computes the
same \emph{function} as the source. That is a statement about one execution at a
time, and confidentiality is not a property of one execution at a time.

Concretely, every one of the following transformations preserves the computed
function and destroys \eqref{eq:ni}:

\begin{itemize}[nosep,leftmargin=1.4em]
\item replacing a constant-time conditional move with a branch, so the observer
      learns the condition from control flow;
\item strength-reducing a division into a multiply-high, changing which
      operations are executed and how long they take;
\item eliminating a store that was scrubbing a secret, because nothing reads it
      afterwards;
\item promoting a secret buffer into registers that are later spilled;
\item choosing an allocation size that depends on a secret count.
\end{itemize}

So a confidentiality checker cannot be a corollary of a correctness checker. It
needs its own property, its own observation model, and --- because the property
is about pairs --- its own machinery.

\subsection{Four outcomes, and why three would be dishonest}

The checker reports one of four values for each question it is asked:

\begin{center}
\begin{tabular}{@{}l p{0.44\linewidth} p{0.27\linewidth}@{}}
\toprule
Outcome & Meaning & Obligations \\
\midrule
\Verified    & no residual forbidden dependence at the claimed boundary
             & must be none \\
\Unsafe      & a forbidden dependence, witnessed by a replayable counterexample
             & must be none \\
\Unknown     & required semantics or a contract is absent
             & must name what is missing \\
\Conditional & acceptable only if named external assumptions hold
             & must name every assumption \\
\bottomrule
\end{tabular}
\end{center}

The two middle-ground values carry the weight. A checker that can only say
\emph{safe} or \emph{unsafe} must guess when it meets a function pointer, a
recursive call, an opaque external callee, an uninitialized read, or a solver
timeout --- and a guess in the \Verified{} direction is exactly the failure that
makes such a tool worthless. \Unknown{} is what lets the analysis fail closed.
\Conditional{} is what lets it be honest about facts that live outside the model
altogether: real helper latency, backend trace preservation, hardware isolation.

\begin{remark}[The invariant that makes the ledger work]
\Verified{} and \Unsafe{} must carry \emph{no} outstanding obligations, while
\Unknown{} and \Conditional{} must name theirs. That is precisely the discipline
a claim ledger encodes: an outcome is only allowed to be terminal when the
ledger for it is empty.
\end{remark}

\subsection{Evidence levels}

The system separates \emph{where} a fact is established from \emph{how much} it
is worth. Five levels are used throughout, and every fixture in the harness
names the ones it relies on.

\begin{center}
\begin{tabular}{@{}l p{0.80\linewidth}@{}}
\toprule
Level & What it supplies \\
\midrule
\Lv{0} & declarations: secrets, observer, release policies, helper contracts,
         hosts, target facts \\
\Lv{1} & propagation: dependence and effects --- branches, addresses,
         latencies, sinks \\
\Lv{2} & a relational witness: two lanes coupled at entry, different secrets,
         different observations \\
\Lv{3} & attribution: relating source, imported IR, and target artifact to a
         compiler boundary \\
\Lv{4} & facts outside the modeled semantics: real latency, compressors,
         hardware isolation, certificates \\
\bottomrule
\end{tabular}
\end{center}

\Lv{0} is a \emph{contract, not a proof of the contract}. If no latency summary
for a helper is available, the honest result is \Unknown{}; selecting a named
operand-dependent test profile lets \Lv{1} classify the call as \Unsafe{} under
that profile, and \Lv{4} is still required before claiming a deployed helper
matches it.

\subsection{How to read the ledger}
\label{sec:ledgerguide}

Each step in Part~\ref{part:walkthrough} carries a panel in the right margin
listing the obligations \emph{still open} after that step. Rows are coloured by
what the step did:

\begin{center}
\begin{tabular}{@{}ll@{}}
\toprule
\textsc{new} (green) & the step opened a new obligation \\
\textsc{reduced} (amber) & the step rewrote an obligation into a different one \\
\textsc{used} (amber) & the step relied on it; the body is unchanged \\
\textsc{discharged} (red) & the step closed it; struck out once, then dropped \\
carried (grey) & untouched by this step, still open \\
\bottomrule
\end{tabular}
\end{center}

Two reading conventions. First, an obligation quantified over an index --- per
coalition, per site, per entry --- is \emph{one} row carrying a subscript, not
one row per instance; otherwise the panel would be unreadable by Step~3.
Second, the bodies are deliberately terse because a struck-out body cannot break
across lines. The abbreviations are:

\begin{center}
\begin{tabular}{@{}ll@{}}
\toprule
$\Sec_\coal(P)$ & $P$ satisfies \eqref{eq:ni} for coalition $\coal$ \\
$\Bind(\artifact,\pipeline)$ & every result is bound to this exact artifact and pipeline \\
$\Total(\manifest)$ & the policy manifest is total over reachable sites \\
$\NFConf(\artifact,\pipeline)$ & the artifact lies inside the analyzable fragment \\
$\Adm \neq \emptyset$ & at least one execution satisfies the declared preconditions \\
$\PairDom_\coal \neq \emptyset$ & at least one admitted low-equal \emph{pair} exists \\
$\DiagOK$ & the diagnostic layer assigned a disposition to every site \\
$\ProdSafe_\coal$ & the exact two-lane product cannot reach a bad state \\
$\ReplayCov_\coal(w)$ & witness $w$ replays under the exact semantics \\
$\PairedRef_\coal$ & the backend preserves the observable trace \\
\bottomrule
\end{tabular}
\end{center}

The arc opens with \eqref{eq:ni} itself and ends with an empty ledger. If you
only read the margins, you will still see which obligation each step traded
away and what it bought.

\section{Background: pairs, products, and 2-safety}
\label{part:background}
\label{sec:twosafety}

This part builds the one piece of machinery the rest of the document leans on.
If you are fluent in information-flow policy but have not thought much about
\emph{how} a relational property gets checked mechanically, this is the gap
worth closing before Part~\ref{part:walkthrough}.

\subsection{Safety properties and their limits}

Most program verification concerns \emph{safety} properties: a predicate over a
single execution, of the form ``no execution ever reaches a bad state.'' Null
dereference, array bounds, assertion violation, and functional correctness
against a specification are all of this shape. The reason this class is
comfortable is that it is checkable by reachability --- build an abstraction of
the state space, show the bad set is unreachable, done. Every mature tool we
have, from abstract interpreters to model checkers to SMT-backed verifiers, is
built for it.

Equation~\eqref{eq:ni} is not of this shape. No single execution of $P$ is a
counterexample to it. A violation is a \emph{pair}: two executions that agree on
everything the observer is entitled to know, and disagree on some observable.
This is the defining feature of a \emph{2-safety} property, also called a
hyperproperty of arity two.

\begin{remark}[Read \eqref{eq:ni} as a specification, not as an encoding]
\label{rem:notencoding}
Equation~\eqref{eq:ni} is stated as a whole-run implication, and it is correct as
an informal statement of intent. It is \textbf{not} how the checker encodes the
question, and turning it directly into a solver query is a known unsoundness.
\cref{sec:ledger} explains why and gives the construction that replaces it.
\end{remark}

\begin{remark}
The practical consequence is worth stating plainly: you cannot instrument $P$,
run it, and check an assertion. There is no assertion over a single run whose
failure means ``leak.'' Dynamic taint tracking, which is what LLVM's
DataFlowSanitizer does, therefore cannot decide \eqref{eq:ni}; it approximates
it by tracking labels through one execution and reporting when a labelled value
reaches a sink. That approximation has no declassification and no notion of
timing, which is why it is the wrong model here even though its label plumbing
is instructive.
\end{remark}

\subsection{The product construction}

There is a standard and very useful move: \textbf{2-safety of $P$ is ordinary
safety of a program that runs two copies of $P$ at once.} Build $P \times P$,
the \emph{product} or \emph{self-composition}, whose state is a pair of $P$
states, one per \emph{lane}. Couple the initial states on what the observer can
see \emph{at entry}, let the secrets range freely, and add a bad-state predicate
that fires when the lanes' observations diverge.

Precisely which facts couple the initial states matters a great deal, and is the
subject of \cref{sec:ledger}. For now note only what they are \emph{not}: they do
not include equality of any release the program has not performed yet.

\begin{center}
\begin{tikzpicture}[
  lane/.style={draw, rounded corners=1.5pt, minimum height=5.5mm, minimum width=14mm,
               font=\small, inner sep=2pt},
  bad/.style={draw, rounded corners=1.5pt, minimum height=5.5mm, minimum width=14mm,
              font=\small, fill=red!12, inner sep=2pt},
  lbl/.style={font=\scriptsize, anchor=center},
  ar/.style={-{Latex[length=1.6mm]}}
]
  \node[lane] (a0) at (0,0.8)   {$\sigma_0$};
  \node[lane] (a1) at (2.2,0.8) {$\sigma_0'$};
  \node[lane] (a2) at (4.4,0.8) {$\sigma_0''$};
  \node[bad]  (a3) at (6.7,0.8) {\textsf{obs}${}=\alpha$};

  \node[lane] (b0) at (0,-0.8)   {$\sigma_1$};
  \node[lane] (b1) at (2.2,-0.8) {$\sigma_1'$};
  \node[lane] (b2) at (4.4,-0.8) {$\sigma_1''$};
  \node[bad]  (b3) at (6.7,-0.8) {\textsf{obs}${}=\beta$};

  \foreach \i/\j in {a0/a1, a1/a2, a2/a3, b0/b1, b1/b2, b2/b3}
    \draw[ar] (\i) -- (\j);

  \foreach \i/\j in {a0/b0, a1/b1, a2/b2}
    \draw[dotted] (\i) -- (\j);

  \node[lbl, anchor=east] at (-0.85,0.8)  {lane 0: $s_0$};
  \node[lbl, anchor=east] at (-0.85,-0.8) {lane 1: $s_1$};

  \node[lbl] at (0,1.7)  {coupled at entry only};
  \node[lbl, text width=2.6cm, align=center] at (6.7,1.85)
       {$\alpha \neq \beta$ is the bad state};
\end{tikzpicture}
\captionof{figure}{The two-lane product. Initial states are coupled on what is
visible \emph{at entry} --- public configuration, public alias topology, and the
initially-visible components --- with the secrets free. Releases are \emph{not}
part of this coupling; they are handled step by step as the lanes run
(\cref{sec:ledger}). Dotted lines are the lockstep alignment. A violation is
exactly a reachable state where the lanes' projected observations differ, which
is an ordinary reachability question over the product.}
\label{fig:product}
\end{center}

This is the hinge of the design. It converts a property no existing tool can
check directly into one that every existing tool can check, at the cost of
squaring the state space. It is also why the specification speaks of an
``audit-all whole-entry product'' rather than a type system: the product is the
\emph{authoritative} object, and a verdict is a reachability result over it.

\subsection{The release ledger, and why whole-run equality is wrong}
\label{sec:ledger}

\begin{remark}[Two things called a ledger]
This document has a \emph{claim ledger} in the margin --- a presentational device
of these notes. The specification separately has a \emph{release ledger}, an
object inside the product semantics. They are unrelated. This subsection is about
the second; every margin panel is the first.
\end{remark}

\cref{rem:notencoding} promised an explanation. Here it is, because this is the
one place where the obvious encoding of \eqref{eq:ni} is not merely imprecise but
\emph{unsound}.

\textbf{What couples the initial states.} The lanes must agree on what the
coalition can see \emph{at entry}: the public configuration, the public alias
topology, and the read-carrier of every initially-visible component. Mechanism
and timing-environment choices are coupled too, so a pair cannot differ on a
modelled implementation choice. Deliberately absent: any requirement that the
secret components differ, and --- the point of this subsection --- any equality of
a release the program has not performed yet.

\textbf{Why that last omission matters.} \eqref{eq:ni} reads: if the authorized
releases come out equal, the observations must come out equal. The obvious
encoding installs $\Rel(p,s_0) = \Rel(p,s_1)$ as an initial constraint on the
product and then asks for bad-state reachability. A two-instruction program shows
that this is wrong:

\begin{center}
\code{emit(public\_channel, h)} \quad;\quad \code{release(policy, audience, h)}
\end{center}

The secret is published at step~1 and only afterwards passed through an authorized
release at step~2. Under the whole-run encoding the constraint keeps only lane
pairs whose complete release histories agree --- and since the released value
\emph{is} $h$, that means only pairs with equal $h$. Those pairs trivially agree at
step~1 too. The leak is never examined: the encoding has filtered away precisely
the pairs that demonstrate it. Note the direction of the error --- it is
permissive, not conservative.

\textbf{The fix is causal.} Releases are consulted step by step as the lanes run,
never installed up front:

\begin{center}
\begin{tabular}{@{}p{0.31\linewidth} p{0.58\linewidth}@{}}
\toprule
At an aligned step & The ledger does \\
\midrule
no release here
  & require the projected words to agree, as usual \\[3pt]
release authorized for this observer, values equal
  & append the entry and continue; the obligation stays active \\[3pt]
release authorized for this observer, values differ
  & permitted --- but everything outside the release's declared footprint must
    still agree, and the pair's obligation then \emph{retires} \\[3pt]
release not authorized for this observer
  & conceal the value; require everything else about the event to agree \\
\bottomrule
\end{tabular}
\end{center}

The transition takes only the current ledger and the current pair of events. It
has no parameter through which a later event can be inspected, so a release that
has not happened cannot excuse an observation that already has.

\begin{remark}
\label{rem:mtcm3}
``Retires'' is doing real work: an authorized difference ends the obligation for
that pair \emph{from that point}, rather than licensing differences
retroactively. A pair is protected up to its first authorized divergence and
released afterwards, which is what makes declassification expressible without
making it a loophole.

This is not hypothetical. The program above is encoded in the harness as
\fx{prefix\_causal\_release.bad.mlir}, and a checker built on the whole-run
encoding reports that fixture safe.
\end{remark}

\subsection{Why the product is harder than it looks}

Figure~\ref{fig:product} is drawn with the lanes in lockstep, which quietly
assumes they follow the same control path. They need not. Once a secret can
influence control flow the lanes desynchronize, and the product must say
something about misaligned executions. This is where the bad-state predicate
stops being ``the outputs differ'' and becomes a disjunction over structural
divergences:

\begin{enumerate}[nosep,leftmargin=1.6em]
\item one lane can step and the other cannot;
\item the two lanes' next control locations differ;
\item their termination or failure statuses differ;
\item their event site or occurrence alignment differs;
\item the projected event words they emit differ.
\end{enumerate}

Items 1--4 are \emph{structural}: they compare the shape of the two executions,
not their data. That distinction matters because structural divergence is not
weakened when the observer's projection conceals a payload. An observer who
cannot read a value can still count operations, notice that a branch was taken,
and see that a run terminated early.

\begin{remark}[Where the subtle bugs live]
Nearly every unsoundness in this area is a case where somebody checked item 5
and forgot items 1--4, or checked a \emph{local} version of item 4 rather than a
global one. Two of the countermodels in Part~\ref{part:walkthrough} are exactly
of this form, and one of them --- \cref{step:product} --- is a program where
every per-site payload and count agrees and only the global ordering differs.
\end{remark}

\subsection{Two layers, and the discipline between them}

Squaring the state space is expensive, so it is tempting to avoid the product
with a cheap syntactic analysis: label each value \textsf{Low} or \textsf{High},
join labels at merges, and complain when \textsf{High} reaches a sink. That is a
classic information-flow type system, and it is genuinely useful. But note
carefully which way its guarantee points. For a \emph{sound over-approximation}
it is \textbf{acceptance} that is conclusive --- the classical theorem is that a
well-typed program \emph{is} noninterferent --- while its \emph{rejections} may
be false alarms. A type system proves; it cannot refute.

Two things nevertheless remove that proving power in this setting. The
observation model here includes timing, addresses, allocation sizes, and
operation counts, and a label discipline over values says nothing about any of
them. And the property is release-relative, so a declassification must be
\emph{permitted} at exactly the right point, which the prefix-causal ledger of
\cref{sec:ledger} decides and a label lattice cannot express. What survives is a
useful \emph{triage}: cheap, sound in the direction of refusal, and unable to
close the question.

The architecture therefore has two layers with a deliberately asymmetric
contract:

\begin{center}
\begin{tabular}{@{}p{0.29\linewidth}p{0.61\linewidth}@{}}
\toprule
\textbf{Diagnostic layer} (\Lv{1}) & one \textsf{Low}/\textsf{High} pass, with
no proof-authoritative strong update, no function summaries, and no slice
selection. Deliberately weak. \\[3pt]
\textbf{Product layer} (\Lv{2}) & the exact two-lane product over the whole
entry. Authoritative. \\
\bottomrule
\end{tabular}
\end{center}

The contract is: \textbf{the diagnostic layer may narrow the work, and may never
discharge an obligation.} Its weakness is the point. Because it cannot perform a
strong update it cannot conclude that a secret store was later overwritten;
because it has no summaries it cannot conclude anything about a callee. So it
cannot be unsound --- it can only be imprecise, and its correct output at a site
it cannot decide is \RelReq{}: \emph{the product must decide this one}.

\begin{remark}[The single most dangerous repair]
\label{rem:danger}
When the diagnostic layer reports a violation on a program that is in fact safe,
the tempting fix is to teach it the missing reasoning --- a strong update, a
summary, a slice. Every such addition moves proof burden into the layer designed
to carry none, and the specification forbids the consequence outright: no
diagnostic disposition, including \emph{statically discharged} or \emph{not
observable}, may establish safety or remove a product constraint. This is why
the harness ships \emph{paired anti-controls} rather than controls alone, a
discipline developed in \cref{sec:pairs}.
\end{remark}

\section{The walkthrough}
\label{part:walkthrough}

Ten steps take \eqref{eq:ni} from an informal wish to an empty ledger. Read the
margin panels as the primary text if you like; the prose exists to justify each
trade.

\clgreset

% ---------------------------------------------------------------- step 1
\begin{step}{What we are actually trying to prove}{%
  \clgnew{NI}{$\forall \coal.\ \Sec_\coal(P)$}%
  \clgnew{BIND}{$\Bind(\artifact,\pipeline)$}%
  \clgnote{BIND}{stays open almost to the end}%
}
\label{step:statement}

Open with the goal and with the one obligation that is easy to forget: a verdict
is a statement about \emph{a particular artifact}. ``This program is secure'' is
not a well-formed claim until you say which bytes you mean.

That second obligation is not pedantry. A confidentiality result computed on
LLVM IR captured before a late pass runs, whose mutated output then proceeds to
instruction selection, is a statement about a module that never executed. The
honest result in that situation is \Unknown{}, and nothing weaker than a binding
to the exact artifact --- bitcode hash, ordered pipeline digest, the capture
point, and the first IR-to-machine-IR boundary --- rules it out.

$\Sec_\coal$ is indexed by a \emph{coalition} $\coal$, a set of principals who
may pool what they observe. \cref{step:coalitions} explains why the index is
there. For now note that a single row in the panel stands for the whole family:
the ledger convention from \cref{sec:ledgerguide} is that an obligation
quantified over an index is one row with a subscript, not one row per instance.
\end{step}

% ---------------------------------------------------------------- step 2
\begin{step}{Declare the policy, and make the declaration total}{%
  \clgtouch{NI}%
  \clgnew{MAN}{$\Total(\manifest)$}%
}
\label{step:manifest}

Nothing about \eqref{eq:ni} is decidable until four questions have declared
answers: which inputs are secret, who observes what, which releases are
authorized and to whom, and which host executes each function.

These are supplied by a \emph{policy manifest}, and the obligation this step
opens is that the manifest is \emph{total} over everything reachable. Partiality
is the failure mode, not incorrectness: a reachable function with no declared
host has no observation projection, so no verdict follows for it --- and the
tempting default of treating an unplaced function as local and trusted is an
invented premise.

Two rules about the manifest are worth stating early because they are the ones
implementations break.

\textbf{Policy may not be inferred from names.} Not from an LLVM value name, not
from an argument position, not from \code{noalias}. This sounds like a style
guideline and is in fact load-bearing: it is what forces the manifest to be a
parsed input rather than a convention. It also has a blunt mechanical
justification --- \code{mlir-opt} discards SSA value names at parse, so a policy
encoded in a name does not survive to the analysis at all.

\textbf{Authority is not one relation.} A single principal-to-clearance map
cannot express the policies people actually want. The manifest instead carries
several independent relations:

\begin{center}
\begin{tabular}{@{}ll@{}}
\toprule
\emph{visibility} & who may \textbf{observe} an item; drives the projection \\
\emph{audience} & who a released value is \textbf{for}; drives the release premise \\
\emph{placement} & which host \textbf{executes} a function \\
\bottomrule
\end{tabular}
\end{center}

Visibility is itself carried by four separate bases --- over hosts, components,
outputs, and errors --- which is the concrete reason a single clearance level per
principal is insufficient: an item's observability is a function of where it
lives, not only of who is asking.

\begin{remark}[A relation the model does not yet have]
\label{rem:authorizers}
A natural fourth relation is \emph{who may authorize a release}, distinct from
who may read the released value. That separation is what would make ``may
declassify but may never read'' expressible --- an auditor trusted to decide
\emph{that} a summary may be published, without being trusted to see the raw
data behind it.

No such relation exists in the specification. Releases are identified by policy
and audience; there is no principal-level authorizer anywhere in the manifest
schema. So the auditor policy above is currently \textbf{not expressible}, and
the design examples that use an \code{authorizers} field are a proposed
extension rather than a description of the system. This is a genuine gap, not a
simplification for exposition.
\end{remark}
\end{step}

% ---------------------------------------------------------------- step 3
\begin{step}{Derive the coalitions; do not accept the authored list}{%
  \clgupd{NI}{$\forall \coal \in \coalitions.\ \Sec_\coal$}%
  \clgnew{CLO}{$\coalitions = \dnclose\coalitions_{\max}$}%
  \clgdone{CLO}%
  \clgnote{CLO}{immediate, and therefore easy to skip}%
}
\label{step:coalitions}

The coalition family is the \emph{downward closure} of the authored maximal
list, and it includes the empty coalition, which represents the world-visible
observer. Every member must be checked, including members nobody wrote down.

This step creates its obligation and discharges it immediately, because the
closure is a finite set operation. The reason to give it a step at all is that
``obvious and immediate'' is not the same as ``done,'' and the failure is silent:
an implementation that iterates the authored list rather than the derived closure
produces confident green results on artifacts that leak.

Two facts make the closure irreducible --- neither is an implementation detail.

\textbf{Verdicts are not monotone in the coalition.} The natural intuition is
that a larger coalition observes more, so checking the largest is the hardest
case and suffices. Release audiences break this. Authorization is
\emph{per-audience}, so whether a given release retires a pair's obligation
(\cref{sec:ledger}) depends on which coalition is asking. For a coalition inside
the declared audience the differing values are permitted and the obligation
retires; for a coalition outside it, nothing retires, the obligation is still
active, and \emph{the identical operation is a leak}. So \Verified{} at one
coalition and \Unsafe{} at another, with no containment between them.

\textbf{Joint visibility is not reconstructible from singletons.} An item may be
declared visible to $\{\textsf{alice},\textsf{bob}\}$ jointly and to neither
alone --- two shares of a secret, recoverable only together. A visibility
relation stored as principal-to-items cannot express this, and a report built by
evaluating singletons and taking unions cannot derive it.

Together these forbid deduplicating rows by coalition monotonicity, and forbid
omitting a derived coalition from the report. The cost is real and structural:
the closure over $n$ principals has $2^n$ members, each needing a row.
\end{step}

% ---------------------------------------------------------------- step 4
\begin{step}{Get inside a fragment you can actually reason about}{%
  \clgnew{NF}{$\NFConf(\artifact,\pipeline)$}%
}
\label{step:normalform}

Real LLVM IR is too large a language to have an exact semantics for. The move is
to fix a closed-world fragment --- a whitelist of types, opcodes, constants, and
intrinsics --- and to define the analysis only there. Everything outside is out
of profile.

The critical rule is what happens on the boundary. If an artifact does not
conform, the result is \Unknown{} with a reason naming the offending construct.
It is emphatically \textbf{not} a counterexample, even when the nonconforming IR
contains what looks unmistakably like a secret-dependent branch. A best-effort
pre-conformance scan may emit a triage finding, but that record is
non-authoritative and may never be presented as a replayable witness.

The reason for that asymmetry is that a counterexample is a claim about the
exact semantics, and outside the fragment there is no exact semantics to make a
claim about. Reporting \Unsafe{} there would be unsound in the same way that
reporting \Verified{} would be --- just less obviously so, and therefore more
dangerous, because a false alarm looks like diligence.

\begin{remark}[A scoping fact worth internalizing]
The fragment is scalar: integers and pointers, plus \code{float} and
\code{double}. Vectors are out of profile, as are the exotic float formats
(\code{half}, \code{bfloat}, and the 80- and 128-bit forms). So an analysis
defined on it says nothing about vector or tensor code until each is explicitly
admitted.

Floating point is admitted but is not thereby free: a reachable arithmetic
operation whose result could be an ambiguous NaN payload is a refusal, because
the fragment does not choose among the payloads the IR permits. Bit-preserving
movement and comparison stay eligible. Either way, a corpus of integer fixtures
does not exercise these refusal paths, and a real workload triggers them
immediately.
\end{remark}
\end{step}

% ---------------------------------------------------------------- step 5
\begin{step}{Check that the question is not vacuous}{%
  \clgnew{VAC}{$\Adm,\PairDom_\coal \neq \emptyset$}%
  \clgnote{VAC}{a claim-adequacy gate, not a premise}%
}
\label{step:vacuity}

Here is a way to obtain a proof of \eqref{eq:ni} for any program whatsoever.
Declare preconditions that contradict each other. There is then no admitted
state, hence no admitted low-equal pair, hence the initial state set of the
product is empty, hence the bad state is unreachable. The formula is
unsatisfiable without executing a single instruction, and the checker reports
success.

This is the countermodel that motivates the gate, and it has a subtler second
form: preconditions that admit exactly one state. Now a nonempty admitted set
and a nonempty pair domain both hold, via the pair of that state with itself,
while no secret component actually \emph{varies}. The product is inhabited and
still proves nothing.

So two things must be exhibited, not assumed: at least one admitted execution,
and at least one admitted \emph{coupled low-equal pair} for each coalition. And a
third disposition is needed for the middle case --- a secret component that
provably cannot vary is reported as constrained-or-unexercised rather than
silently contributing to a proof, and the verifier may not infer an
expected-to-vary assertion from a component's name, width, or classification.

\begin{remark}
This gate is not a logical premise of \eqref{eq:ni}; it is a
\emph{claim-adequacy} gate on the report. The distinction matters for the data
structure: both kinds of blocker produce \Unknown{} with a reason, but only
premises appear in the antecedent of the theorem being claimed, so a report that
cannot tell them apart cannot state which theorem it proved.
\end{remark}
\end{step}

% ---------------------------------------------------------------- step 6
\begin{step}{Run the diagnostic layer, and let it fail}{%
  \clgnew{DIAG}{$\DiagOK$}%
  \clgtouch{MAN}%
}
\label{step:diag}

Now the cheap pass from \cref{sec:twosafety}: one \textsf{Low}/\textsf{High}
propagation assigning every site a disposition. Its obligation is only that the
assignment is \emph{total} --- every site classified, nothing silently skipped.

The disposition domain is where the discipline lives:

\begin{center}
\begin{tabular}{@{}ll@{}}
\toprule
not observable & no coalition-visible payload claim at this site \\
statically discharged & closed by a named rule, with premise references \\
definite violation & a violation with a witness recipe \\
\RelReq{} & precision lost; \textbf{the product must decide} \\
\Unknown{} & this \emph{analysis} did not complete: a health failure \\
\bottomrule
\end{tabular}
\end{center}

The last two rows are not two grades of the same thing, and the corpus records
their conflation as a resolved defect. \RelReq{} is a \emph{precision} outcome:
it is this layer working exactly as designed. The diagnostic \Unknown{} is a
\emph{health} failure --- incomplete traversal, a malformed record, a failed
consistency check --- and it is the only diagnostic disposition that blocks
\Proved{}. Keeping them apart is what lets a run say ``I was imprecise here''
without also saying ``I may be broken.''

The last row is narrower than it looks and is worth not misreading. A
\emph{diagnostic} \Unknown{} is reserved for a failure of the diagnostic pass
itself --- it did not achieve its required whole-entry coverage. That is a
different thing from the artifact-level \Unknown{} of \cref{step:normalform} and
\cref{step:verdict}, which reports missing semantics or an absent contract. Two
layers, two \Unknown{}s, and only the second is a verdict about the program.

The fourth value is the one that makes the design work, and the one an
implementer is most likely to try to eliminate. Consider four programs that are
genuinely safe and that this layer cannot see are safe:

\begin{itemize}[nosep,leftmargin=1.4em]
\item a secret written to a slot and then fully overwritten by a public value
      before any read --- needs a strong update, which this layer may not do;
\item a secret at one byte offset with the sink fed from a disjoint offset ---
      needs byte-exact disjointness, and the diagnostic's memory regions are not
      proof certificates;
\item \code{xor} of a value with itself --- needs value congruence, which a
      label lattice collapses away;
\item a branch on a secret whose two edges target the \emph{same} block ---
      structurally lockstep, so disjunct 2 of \cref{sec:twosafety} cannot fire.
\end{itemize}

All four are \RelReq{} here and \Proved{} at the product. None of them is
``silence''; demanding silence is what pushes an implementer toward
\emph{statically discharged} or \emph{not observable}, which \cref{rem:danger}
forbids. There is even a reserved reason code for having taken that shortcut.

The fourth example carries a trap worth its own sentence. The obvious repair ---
``identical successors imply no control leak'' --- is unsound, and
\cref{step:product} exhibits the program it accepts.
\end{step}

% ---------------------------------------------------------------- step 7
\begin{step}{Build the exact product}{%
  \clgnew{PROD}{$\ProdSafe_\coal$}%
  \clgdone{DIAG}%
  \clgtouch{NF}%
}
\label{step:product}

Construct the two-lane product over the whole entry, with the bad-state
disjunction from Part~\ref{part:background}, and ask whether a bad state is
reachable. The diagnostic layer's obligation discharges here: its job was triage,
and the product now carries the burden.

Line~\ref{ln:rel} is the exception that makes the construction correct rather
than merely strict. A release event does not go through the disjuncts at all; it
is routed to the prefix-causal transition of \cref{sec:ledger}, where an
\emph{authorized} release whose two values \emph{differ} is permitted and retires
the pair's obligation. Without that routing every authorized declassification
would register as a leak, and the checker would reject every useful program.

\begin{algorithm}[H]
\caption{Bad-state predicate for coalition $\coal$ at aligned step $k$. Applies
\emph{outside} the release transition of \cref{sec:ledger}; a release event is
routed there instead, and an authorized release with differing values is
permitted rather than bad.}
\label{alg:bad}
\begin{algorithmic}[1]
\State \label{ln:rel} \textbf{if} this step is a release event \textbf{then return} \textsc{ledgerStep}
\State \label{ln:step} \textbf{if} exactly one lane can take a step \textbf{then return} \textsc{bad}
\State \label{ln:loc}  \textbf{if} the lanes' next control locations differ \textbf{then return} \textsc{bad}
\State \label{ln:status} \textbf{if} their termination or failure classes differ \textbf{then return} \textsc{bad}
\State \label{ln:align} \textbf{if} the emitted event's site or occurrence index differs \textbf{then return} \textsc{bad}
\State \label{ln:word} \textbf{if} $\Proj_\coal(\text{event}_0) \neq \Proj_\coal(\text{event}_1)$ \textbf{then return} \textsc{bad}
\State \textbf{return} \textsc{ok}
\end{algorithmic}
\end{algorithm}

Three places where an independent implementation plausibly diverges, and where
divergence is unsound rather than merely different:

\textbf{Line~\ref{ln:loc} keys on successor identity, not on the condition's
label.} A branch on a secret into one block has a singleton successor set and
cannot fire this line, which is exactly why the fourth example in
\cref{step:diag} is safe. But now consider a branch on a secret into one block
whose two \emph{edges carry different operands} --- in the LLVM dialect a
$\phi$ is a block argument, so the merged value arrives as an argument of the
target block with no operand edge from the branch. Lines~\ref{ln:step}
through~\ref{ln:align} all pass. The leak lands only on Line~\ref{ln:word}.

That is the program the repair rejected in \cref{step:diag} would accept, and it
is not hypothetical: canonicalizing an ordinary two-armed diamond in
\code{mlir-opt} produces precisely this shape. It also refutes a more general
principle --- that an ordinary SSA slice is closed around a $\phi$ --- since a
dependence relation built over operand edges never sees the gating fact.

\textbf{Line~\ref{ln:align} is global, not per-site.} Comparing occurrence counts
\emph{per event template} is not the same as comparing the global interleaving.
Two lanes can emit the same events, the same number of times, with identical
payloads, in the opposite order. Every per-template row agrees; the traces
differ. So the projected trace must be an ordered sequence with pairwise index
alignment and a first-divergence position --- never a multiset, and never a
map from site to last value.

\textbf{Lines~\ref{ln:step}--\ref{ln:align} do not weaken under projection.}
They are structural, so a coalition that cannot read a payload still
distinguishes these executions. An implementation that applies the coalition
projection before evaluating them will wrongly equate lanes that differ in
control location, count, or termination.

\begin{remark}[What may not be filtered]
An execution that exhausts an unrolling bound, faults, or risks undefined
behaviour is \emph{retained}, not deleted. Deleting it silently narrows the proof
domain --- the same vacuity failure as \cref{step:vacuity} arriving by a
different route. Bound adequacy is a universal obligation discharged separately,
not a filter applied inside the admitted set.
\end{remark}
\end{step}

% ---------------------------------------------------------------- step 8
\begin{step}{Prove, refute, or refuse}{%
  \clgdone{PROD}%
  \clgdone{VAC}%
  \clgnew{WIT}{$\ReplayCov_\coal(w)$}%
  \clgdone{WIT}%
}
\label{step:verdict}

The product query returns one of three things, and the third is not a failure of
the tool.

\textbf{Unreachable} $\Rightarrow$ \Proved{}, provided the adequacy gate of
\cref{step:vacuity} holds. That proviso is why both obligations discharge in the
same step.

\textbf{Reachable, with a witness} $\Rightarrow$ \Ctrex{}, but only if the
witness \emph{replays}. A model returned by a solver is a candidate, not a
finding: it must be re-executed under the exact semantics and shown to reach the
bad state, with its prefix covered. Failed replay or missing coverage is
\Unknown{} and tool inconsistency --- never a counterexample. This is the
obligation that opens and closes inside this step.

\textbf{Neither} $\Rightarrow$ \Unknown{}, with the blocker set named.

Two rules about aggregation, both of which are easy to get wrong and neither of
which is cosmetic.

\emph{A replayed witness outranks an open universal obligation.} If a witness
reaches a bad state and replays, the verdict is \Ctrex{} regardless of an
unfinished bound-adequacy proof, an unfinished definedness proof, or another
solver timeout elsewhere --- and regardless of whether a solver, a fuzzer, a
differential search, or a person found it.

Source-independence is a requirement on the \emph{verdict}, not on the record.
The result must be identical from every discovery route, and no step of the
argument may consult provenance. The record itself is the opposite of anonymous:
it must \textbf{retain} the discovery source and the seed or model, precisely so
the candidate can be re-executed and the replay audited. Confusing the two ---
stripping provenance in the name of source-independence --- destroys
reproducibility while buying nothing, since it is the argument that must ignore
the source, not the file.

\emph{Otherwise every blocker is preserved.} With one blocker the result is
\Unknown{} carrying it. With several it is \Unknown{} carrying a compound reason
over the \emph{complete ordered} blocker set. A flat enumeration of reason codes
cannot represent that, which makes the compound constructor a genuine
requirement on the record type rather than a convenience.

\begin{remark}
Note the shape of the record this forces. \Unsafe{} carries a replayable
witness; \Unknown{} carries an ordered blocker set; \Proved{} carries neither.
A single ``reason'' string field cannot hold all three, and a design that tries
will end up attaching reason codes to proofs, which are reason-free.
\end{remark}
\end{step}

% ---------------------------------------------------------------- step 9
\begin{step}{Close the deployment gap, or admit it is open}{%
  \clgnew{DEP}{$\PairedRef_\coal$}%
  \clgnote{DEP}{leaving this open is what \Conditional{} means}%
}
\label{step:deployment}

Everything so far concerns a frozen artifact under a declared machine model. Two
gaps remain between that and a running system, and they are of different kinds.

\Lv{3} is \emph{attribution}: relating source, imported IR, and the final target
artifact, so a violation or a repair can be ascribed to a specific compiler
boundary. This is where the classic ``the compiler broke my constant-time code''
findings live --- a conditional move lowered to a branch on a target without
conditional moves, or a division strength-reduced into a multiply-high so the
operation you were reasoning about is not the operation that executes.

\Lv{4} is \emph{facts outside the modeled semantics}: real helper latency,
compressor behaviour, hardware isolation between tenants, sanitizer sufficiency,
certificate soundness, and whether the backend preserved the observable trace at
all. These are not provable from the IR. They are empirical or hardware facts.

The honest handling is the fourth outcome. A repair that is correct in the model
but depends on an undischarged external fact is \Conditional{}, and must name
every assumption it rests on. An \Unsafe{} or \Verified{} result may name none.

\begin{remark}[Why this obligation is per-coalition and separate]
Deployment status is an independent axis from model status. A model-level
counterexample cannot be repaired by deployment evidence, and deployment evidence
cannot upgrade an \Unknown{} model result. Collapsing the two into one verdict
token is the most common structural error, and it shows up concretely: an
obligation like ``the backend preserved the trace'' and one like ``this helper's
latency is operand-independent'' belong to different axes and should not share a
field.
\end{remark}
\end{step}

% ---------------------------------------------------------------- step 10
\begin{step}{Close the arc}{%
  \clgdone{DEP}%
  \clgdone{NF}%
  \clgdone{MAN}%
  \clgdone{BIND}%
  \clgdone{NI}%
}
\label{step:close}

Discharge the artifact-level gates and the goal. The manifest was total, the
artifact conformed, the results are bound to those exact bytes and that exact
pipeline, and the deployment obligations of \cref{step:deployment} were
\emph{discharged}. The product was unreachable under a non-vacuous admitted
domain, so \eqref{eq:ni} holds for every derived coalition at the claimed
boundary.

This is the \Verified{} branch, and it is the only branch in which the arc closes.
Had the deployment obligations merely been \emph{named} rather than discharged,
the outcome would be \Conditional{} --- and the ledger would end with that row
still open, which is exactly the invariant from Part~\ref{part:orientation}:
\Verified{} may carry no obligations, so a non-empty ledger is not a proof but a
different verdict. An empty ledger is the picture of one specific outcome, not of
the analysis having run.

What has \emph{not} been established is worth saying explicitly, because it is
the difference between an honest result and an overclaim:

\begin{itemize}[nosep,leftmargin=1.4em]
\item nothing about constructs outside the fragment of \cref{step:normalform};
\item nothing about a different artifact, pipeline, or target profile;
\item nothing about real latency, hardware isolation, or speculation beyond what
      \cref{step:deployment} explicitly named and discharged;
\item nothing about coalitions of principals absent from the manifest, since the
      closure was taken over the declared set.
\end{itemize}

Every one of those is a boundary, not a bug. The value of the four-outcome
contract is that it can state them.
\end{step}

\clgboard

\noindent
The ledger is empty. Most middle steps traded an obligation for a successor that
was easier to attack; a few only opened one, because some preconditions have to
be stated before anything can be traded against them. The two that stayed open
longest were the least glamorous: that the policy was total, and that the answer
was about the artifact we actually shipped.

\section{What the harness actually establishes}
\label{part:harness}

Part~\ref{part:walkthrough} describes a checker that does not exist yet. What
exists is a regression corpus that pins the obligations the checker will have to
discharge. This part explains what that corpus does and --- more usefully --- what
it does not.

\subsection{Two kinds of test}

The suite contains 45 IR fixtures and 7 integration tests, and they split into
two categories that behave completely differently.

\textbf{Green tests are facts locked in.} They assert something true right now
about the IR, the metadata, or the toolchain. They catch regressions.

\textbf{Red tests are executable specifications.} Seventeen fixtures fail on
purpose. Each carries a diagnostic run with an expected diagnostic naming the
exact stable reason a future analysis must emit at a specific operation, and each
fails with \emph{expected error not produced}. They are bring-up gates, not
disabled tests: implementing the analysis and inserting it into those runs is
what turns them green. A suite where every test passes would be a suite that had
stopped specifying anything.

\subsection{Four mechanical layers}

\begin{center}
\begin{tabular}{@{}l p{0.62\linewidth}@{}}
\toprule
Layer & What it establishes \\
\midrule
IR validity     & every fixture parses and satisfies dialect invariants \\
Shape           & the decisive operation --- store, branch, address, division,
                  allocation, backedge --- is present and connected \\
Shape stability & for the eight fixtures carrying a \code{--canonicalize} run
                  --- the four precision controls, the allocation pair, the loop,
                  and the predecessor-choice anti-control --- the decisive shape
                  survives canonicalization, so those cannot silently decay into
                  a different scenario. The other 37 pin shape only in the form
                  they were written \\
Metadata        & each fixture has exactly one of the four outcomes, one
                  observer, a stable reason, obligations required precisely when
                  the outcome is \Unknown{} or \Conditional{}, and real
                  provenance \\
\bottomrule
\end{tabular}
\end{center}

Note what is absent. \textbf{No layer checks an information-flow result},
because no analysis is wired in. The corpus states the specification; it does not
yet test an implementation of it. Being clear about this is the difference between
a useful artifact and a misleading one.

\subsection{Controls, and why every control needs an anti-control}
\label{sec:pairs}

A corpus of leaky programs is satisfied by an analysis that reports every program
as unsafe. That is not a hypothetical failure --- it is the cheapest way to pass a
suite built only from violations, and it makes the tool worthless without making
any test go red.

The guard is a family of \emph{precision controls}: fixtures that are genuinely
noninterferent, where a reported violation is imprecision rather than a finding.
The four in the corpus are exactly the four examples from \cref{step:diag}:
identical-successor branch, value cancellation, public overwrite before
observation, and offset-disjoint reload.

But controls alone are unsafe to ship, because the cheapest way to quiet a
control is usually an unsound repair. So a control wants a paired
\textbf{anti-control} that must retain the opposite outcome:

\begin{center}
\begin{tabular}{@{}p{0.34\linewidth} p{0.54\linewidth}@{}}
\toprule
Control & Unsound repair its partner blocks \\
\midrule
identical-successor branch
  & ``identical successors imply no control leak'' --- the partner is the same
    single-successor branch with differing block-argument operands \\[3pt]
public overwrite
  & adding a proof-authoritative strong update to the diagnostic layer \\[3pt]
offset-disjoint reload
  & coarsening offsets to a cache line, which is a refusal and never a proof \\[3pt]
public world-structural allocation size
  & refusing every dynamic allocation, or treating an equal public cap as an
    equal size \\
\bottomrule
\end{tabular}
\end{center}

The rule is: \emph{never satisfy one member of a pair by weakening the other.}
The first row is the sharpest instance, and it is not a stylistic worry. Running
\code{mlir-opt}'s canonicalizer on an ordinary two-armed diamond produces a
single-successor branch whose edges carry different operands --- byte-identical in
shape to the control, differing only in the operands. An implementation repaired
by that rule accepts a real leak, silently.

\subsection{Which step each fixture guards}

\begin{center}
\begin{tabular}{@{}l p{0.58\linewidth}@{}}
\toprule
Step & Fixtures \\
\midrule
\cref{step:manifest}    & the wrong-party and wrong-host families; the actor
                          examples for the four relations \\
\cref{step:coalitions}  & the downward-closure example, where the leak exists
                          only at a coalition nobody authored \\
\cref{step:normalform}  & currently \emph{uncovered}: no fixture exercises a
                          nonconformance refusal \\
\cref{step:vacuity}     & currently \emph{uncovered}: the vacuity countermodel
                          is not encoded \\
\cref{step:diag}        & the four precision controls \\
\cref{step:product}     & the predecessor-choice anti-control; trace order and
                          multiplicity are \emph{not} yet encoded \\
\cref{step:verdict}     & the unproved-alias and bound-exhaustion refusals \\
\cref{step:deployment}  & the CKKS and wolfSSL helper-latency families, the only
                          \Conditional{} fixtures \\
\bottomrule
\end{tabular}
\end{center}

\subsection{One call, four outcomes}

The wolfSSL 3579 family is the most instructive group in the corpus, because it
demonstrates that the observer model is a load-bearing \emph{input} rather than a
label. \emph{Three} fixtures share one byte-identical body --- a wide multiply
lowered to \fx{@\_\_muldi3} on a 32-bit target --- and differ only in the
declared helper profile. Two more complete the family without sharing that body:
the un-lowered source, whose whole body is a single \code{llvm.mul}, and a
repaired constant-time replacement, which is a 64-iteration mask/add loop that
emits no helper call at all:

\begin{center}
\begin{tabular}{@{}l l p{0.38\linewidth}@{}}
\toprule
Variant & Outcome & Declared knowledge \\
\midrule
source                   & \Unknown{}     & no target lowering, no helper contract \\
target, no contract      & \Unknown{}     & the call exists; no latency summary \\
target, affected profile & \Unsafe{}      & a named operand-dependent helper profile \\
target, constant latency & \Verified{}    & a named constant-latency profile \\
target, repaired         & \Conditional{} & repaired, pending target evidence \\
\bottomrule
\end{tabular}
\end{center}

Three of the four outcome values --- \Unknown{}, \Unsafe{}, \Verified{} --- on
one byte-identical body, driven entirely by declarations. The \Conditional{} row
is different in kind: it comes from a code repair, not a declaration change, and
saying otherwise would overstate an already strong point. That
is the cleanest demonstration available of the point in \cref{step:manifest}:
\Lv{0} is a contract, not proof of the contract, and the same instruction is
\Unknown{}, \Unsafe{}, or \Verified{} according to what has been declared about
the machine it runs on. Note also that the two \Unknown{} rows are genuinely
different scenarios with different missing facts, not one file with an
undecided verdict.

\subsection{Countermodels as regression tests}

Four of the metatheory's required negative results are encoded. These are the
tests a \emph{re-implementation} would fail: they capture reasoning errors rather
than coding errors, and reasoning errors are the ones that survive review.

\begin{center}
\begin{tabular}{@{}p{0.26\linewidth} l p{0.42\linewidth}@{}}
\toprule
Fixture & Outcome & Invalid principle refuted \\
\midrule
predecessor choice     & \Unsafe{}  & an SSA slice is closed around a $\phi$ \\
prefix-causal release  & \Unsafe{}  & a future release may condition an earlier observation \\
alias unproved         & \Unknown{} & unproved ABI separation may be assumed \\
bound exhausted        & \Unknown{} & bounded-run filtering is a sound proof domain \\
\bottomrule
\end{tabular}
\end{center}

The two \Unknown{} rows deserve their outcome explained, because \Unknown{} reads
as a weaker result than it is. Unproved aliasing yields neither safety nor a
counterexample: with no proved disjointness clause and no declared may-alias
realization in the product, asserting \Unsafe{} would itself be unsound, since a
counterexample requires a witness that replays under the exact semantics.
Likewise a reachable bound-exhausted transition that yields no replayed bad
execution is \Unknown{} with the bound-adequacy obligation open --- promoting it
to a counterexample is exactly what \cref{step:verdict} forbids.

The bound-exhaustion fixture is the corpus's only loop with a
\emph{secret-dependent} trip count. Two earlier fixtures already contain
fixed-bound loops with loop-carried block arguments ---
\fx{secret\_embedding\_index.fixed}, a 16-iteration public-induction table scan,
and \fx{wolfssl\_3579\_mul.target\_fixed}, a 64-iteration mask/add multiply ---
so backedge joins were exercised before it. What was absent is a backedge whose
iteration count depends on a secret, and therefore any exercise of bound
adequacy at all.

\subsection{The honest gaps}

The outcome distribution is \Verified{} 21, \Unsafe{} 17, \Unknown{} 5,
\Conditional{} 2. For a tool whose central claim is graceful degradation, five
refusals is thin: there is no fixture for an indirect call, recursion, inline
assembly, a volatile or atomic access, an uninitialized read, or ambiguous
floating-point. Each of those is a place the analysis will meet something it
cannot model, and an untested refusal path is where a silent \Verified{} comes
from.

Three further gaps are worth naming precisely:

\begin{itemize}[nosep,leftmargin=1.4em]
\item About twenty of the older fixtures encode their entire safe/unsafe
      distinction in SSA value names, which \code{mlir-opt} discards at parse.
      After parsing, such a fixture is two indistinguishable stores to two
      unannotated pointers --- so its bring-up gate is not merely unimplemented,
      it is unsatisfiable in principle. The newer fixtures carry policy in
      discardable attributes, which do survive parsing and canonicalization.
\item The corpus is entirely scalar, with three loops, only one of which has a
      non-constant trip count, so nothing exercises
      fixpoint termination at scale, and nothing exercises the vector,
      floating-point, or tensor refusals that any real workload triggers
      immediately.
\item The precision controls carry no diagnostic run, deliberately: their correct
      disposition is \RelReq{}, and a zero-directive run would encode the wrong
      requirement. So they pin shape today and become live precision gates only
      once the record can name an expected disposition. Concretely, the suite
      can currently catch an analysis that \emph{misses} a leak but not one that
      \emph{invents} one.
\end{itemize}

\section{How would you know the checker is correct?}
\label{part:correctness}

``Is it correct?'' is two questions wearing one coat, and separating them is most
of the work.

\begin{enumerate}[leftmargin=1.6em]
\item \textbf{Is the specification right?} Does \eqref{eq:ni} capture what you
      mean by confidentiality? This is \emph{not} provable. It is a definition.
\item \textbf{Is the implementation right with respect to the specification?}
      This is provable, in layers, at wildly varying cost.
\end{enumerate}

Conflating them produces the characteristic formal-methods overclaim: a
mechanized proof of question 2 presented as if it settled question 1.

\subsection{Validating the specification: countermodels, not proofs}

You cannot prove a definition. What you can do is show that plausible
\emph{alternative} definitions are wrong, which is what a countermodel is for: a
small executable program witnessing that some appealing principle fails.

The metatheory names seven such required negative results, and they are worth
reading as a validation \emph{method} rather than a list. Each is a principle
somebody would naturally adopt, plus a program refuting it:

\begin{center}
\begin{tabular}{@{}p{0.52\linewidth} p{0.38\linewidth}@{}}
\toprule
Appealing but invalid principle & Refuted by \\
\midrule
bounded-run filtering is a sound proof domain
  & a secret selecting a trip count past the bound \\[2pt]
a future release may condition an earlier observation
  & emit-then-release, two operations \\[2pt]
local per-template checks determine global trace order
  & two events, same payloads, opposite order \\[2pt]
an ordinary SSA slice is closed around a $\phi$
  & the predecessor-choice diamond \\[2pt]
unproved ABI alias separation may be assumed
  & store through $p$, load through $q$, $p=q$ admitted \\[2pt]
an empty premise domain supports a meaningful proof
  & contradictory preconditions \\[2pt]
unary refinement preserves confidentiality
  & a refinement safe alone, unsafe when paired \\
\bottomrule
\end{tabular}
\end{center}

Four of these are encoded in the harness today. The value of encoding them is
that they are the tests a \emph{re-implementation} would fail --- they capture the
reasoning errors, not the coding errors, and reasoning errors are the ones that
survive code review.

\begin{remark}
This is the strongest available answer to ``how do I know the property is the
right property.'' It is an argument from the absence of known counterexamples to
the definition, strengthened by having actively searched for them. That is
weaker than a proof and much stronger than an assertion.
\end{remark}

\subsection{A ladder of assurance for the implementation}

Order these by cost per unit of confidence, cheapest first. The interesting
observation is that the cheapest three are rarely done, and they subsume most of
what people hope to get from the expensive one.

\subsubsection*{Rung 1: metamorphic and differential testing}

Properties of the checker itself, requiring no ground truth:

\begin{itemize}[nosep,leftmargin=1.4em]
\item \textbf{determinism} --- two runs on one input produce byte-identical
      reports. Worth a deliberate worklist-shuffling switch so the guarantee is
      falsifiable rather than accidental; every pointer-keyed hash map on the
      report path is a hazard here.
\item \textbf{permutation invariance} --- renaming SSA values or reordering
      independent operations must not change the verdict. Note this splits in
      two: the \emph{verdict} must be invariant, while the serialized blocker
      sequence must be \emph{canonical}. A single test comparing whole reports
      conflates them.
\item \textbf{monotone refusal} --- obfuscating the input may turn \Verified{}
      into \Unknown{}, never the reverse.
\item \textbf{idempotence} --- re-running on the checker's own normalized output
      agrees.
\end{itemize}

\subsubsection*{Rung 2: an exhaustive small-domain oracle}

This is the highest-value step available and the one most often skipped.

For \emph{small} bit-widths the property is decidable by brute force. Take a
generated program over, say, 4-bit integers and a handful of memory cells.
Enumerate every pair of secret inputs whose entry states are coupled --- equal
public configuration, equal public alias topology, equal initially-visible
components --- run both lanes concretely, and compare the projected traces step by
step, routing release events through the prefix-causal transition of
\cref{sec:ledger}. That yields \textbf{ground truth} --- not an approximation ---
for that program.

\begin{remark}[The oracle must not use the whole-run encoding either]
It is tempting to write the enumeration the short way: keep only the secret pairs
whose end-of-run authorized releases agree, then compare traces. That is
\cref{rem:mtcm3} all over again, and it makes the oracle agree with a buggy
checker rather than catch it --- including on the harness's own
\fx{prefix\_causal\_release.bad.mlir}. An oracle built on the refuted encoding
validates nothing. Couple at entry; retire obligations as releases occur.
\end{remark}

Now you have a soundness oracle, and you can fuzz against it:

\begin{itemize}[nosep,leftmargin=1.4em]
\item generate random programs in the fragment plus random manifests;
\item compute the brute-force verdict and the checker's verdict;
\item any program where brute force finds a leak while the checker reports
      \Verified{} is \textbf{a soundness bug}, reduced automatically to a minimal
      fixture;
\item the reverse direction is not a bug but a precision measurement, and
      tracking its rate over time is the only honest way to report precision.
\end{itemize}

Two things make this unusually effective here. First, the analyzable fragment is
a \emph{closed-world whitelist}, so a generator covering it is a finite,
enumerable engineering task rather than an open-ended one. Second, every
soundness bug it finds arrives as a concrete, minimal, checked-in regression
fixture --- exactly the artifact the corpus wants. A mechanized proof, by
contrast, tells you a theorem holds but hands you no fixture.

\subsubsection*{Rung 2b: cross-level differential}

Rung 2 fuzzes one level against ground truth. The cheaper sibling fuzzes one
level against \emph{another level of the same program}: analyse before and after
each lowering and compare verdicts. No oracle is needed, because the two verdicts
are each other's oracle.

\Verified{} above and \Unsafe{} below is laundering --- the lowering introduced
the leak. The reverse is evidence destruction. Degradation to \Unknown{} is
fine. This is the check that catches the entire class in
Part~\ref{part:lowering}, it reuses the corpus unchanged, and it is the reason
that part argues for extraction at several levels rather than one.

\subsubsection*{Rung 3: certificate-producing analysis}

Rather than proving the analyzer correct, have it emit a \emph{certificate} that
an independent, much smaller checker validates: an unsatisfiability proof from the
solver, the product encoding, and the mapping from IR sites to product
constraints.

The payoff is a shrunken trusted computing base. You stop needing to trust a
large, fast-changing analysis and start needing to trust a small checker plus the
semantics. This is the same trade that makes proof-carrying code and validated
translation attractive, and it has a decisive practical advantage: the analysis
can be rewritten freely without invalidating anything, because each run
re-establishes its own validity.

\subsubsection*{Rung 4: mechanized proof of the core}

Worth doing --- for a carefully chosen core, discussed next.

\subsection{Mechanizing it for MLIR: what is and is not well-posed}

Start with the uncomfortable part. \textbf{``Prove this correct for MLIR'' is not
a well-posed goal.} MLIR is a meta-IR: an extensible framework whose dialects are
open and user-defined. There is no fixed language to have a semantics for, so
there is nothing to quantify over. Any mechanization must first fix a dialect
subset and give \emph{that} a semantics.

The good news is that this project has already done that, without framing it as a
prerequisite for proof. The normal form of \cref{step:normalform} \emph{is} the
fixed subset: a closed-world whitelist of types, opcodes, constants, and
intrinsics. That is the object to mechanize.

\begin{remark}[Reusable prior work, checked]
\label{rem:priorwork}
The foundation is in better shape than the framing above suggests, and it points
at LLVM-IR level rather than MLIR level.

\emph{Vellvm} is a mechanized semantics for a large sequential subset of LLVM IR,
actively maintained against Rocq, and since its 2021 rewrite it is structured
around interaction trees rather than the original 2012 relational semantics.
Decisively for this programme, \emph{noninterference has already been given a
semantics in the same interaction-tree idiom}, by overlapping authors. That means
the two halves --- an LLVM semantics and a noninterference semantics --- exist in
one compatible formalism, which is the single largest reuse opportunity
available and a strong argument for expressing the fragment at LLVM-IR level.

\emph{Alive2} does bounded translation validation of LLVM optimizations via SMT.
Not a proof, but exactly the right shape for validating the normalizer, per
milestone~6 of the ladder below. (``Milestone'' rather than ``step'' throughout
this part: \emph{Step} is reserved for the numbered steps of
Part~\ref{part:walkthrough}.)
\emph{CompCert} proves functional correctness of a C compiler: architecturally
instructive, different property. Its constant-time-preserving successor is the
closer analogue and is discussed in \cref{sec:positioning}.

For MLIR specifically, mechanized semantics exist in Lean and, separately, as
dialects carrying semantics lowered to SMT --- the latter has already found
miscompilations upstream. Neither carries an information-flow theory.
\end{remark}

\subsubsection*{The decomposition that makes it tractable}

Split the goal three ways and give each part the treatment it can actually
support:

\begin{center}
\begin{tabular}{@{}l p{0.42\linewidth} p{0.28\linewidth}@{}}
\toprule
Part & Content & Treatment \\
\midrule
(a) & \eqref{eq:ni} is the right property & definition; reviewed and
      countermodel-validated, \emph{not} proved \\[3pt]
(b) & the product construction implies the property & mechanized once; this is
      the core theorem \\[3pt]
(c) & this run's artifacts are valid & per-run certificate checking, not proof \\
\bottomrule
\end{tabular}
\end{center}

Almost all the leverage is in refusing to prove (c). The analysis, the
normalizer, and the solver integration will change constantly; proving them once
means re-proving them forever. Validating each run costs a fixed engineering
investment and survives arbitrary rewrites.

\subsubsection*{The core theorem, stated}

The load-bearing obligation is:

\begin{theorem}[Product soundness]
\label{thm:core}
For every artifact $\artifact$ in the fragment, every entry $\entry$, and every
coalition $\coal$: if no bad state of the two-lane product is reachable from the
coupled initial states, then $\Sec_\coal$ holds at $\entry$ --- that is,
$\ProdSafe_\coal(\artifact,\entry) \Rightarrow \Sec_{\coal,\entry}(\artifact)$.
\end{theorem}

Everything else the checker does is a strategy for deciding the antecedent. This
direction is already proved: it is Theorem R8.3 of the metatheory, with the full
premise chain assembled in R9.1 and the solver form in R9.2 --- an accepted
\textsc{unsat} for the identity-bound bad-reachability formula establishes
$\Sec$, and \textsc{unknown} does not.

\subsubsection*{The completeness direction is available, and unstated}

It is tempting to say that only soundness is needed and that completeness is
deliberately abandoned in exchange for the right to answer \Unknown{}. That is
not quite the situation, and the difference is worth a theorem.

The metatheory's Theorem R6.3, \emph{reachable-bad correspondence}, is a
\textbf{biconditional}: a bad state is reachable in the product \emph{if and only
if} there exist admitted states and coupled choices, satisfying low-equivalence,
whose unary executions have a first mismatch forbidden by the observation
relation. The $\Leftarrow$ direction is the soundness of
\cref{thm:core}. The $\Rightarrow$ direction, contraposed, is precisely
$\Sec_\coal \Rightarrow \ProdSafe_\coal$ --- the completeness direction --- and
the corpus never states it as a theorem, corollary, or remark.

One qualifier is load-bearing rather than decorative. R6.3 quantifies over a
\emph{first mismatch forbidden by the active prefix-causal relation}, not over
pairs pre-filtered by whole-run release equality. So the corollary below holds
for $\Sec_\coal$ read prefix-causally, as constructed in \cref{sec:ledger}, and
is \textbf{false} under the naive whole-run reading of \eqref{eq:ni}: the
\code{emit-then-release} program satisfies the whole-run biconditional while
leaking. Stating the completeness result without that qualifier would convert a
correct theorem into an unsound one.

\begin{theorem}[Relative completeness, harvestable from R6.3]
\label{thm:complete}
For an artifact in the fragment, with every contract present, within the declared
execution bound, and relative to the fixed observation relation:
$\ProdSafe_\coal \Leftrightarrow \Sec_\coal$.
\end{theorem}

Read that carefully, because the qualifiers carry it. It says the product is
\emph{exact} within the bound: no false negatives, and --- the part a type system
can never claim --- no false positives either. It does \emph{not} say the
observation relation captures every physical channel; the metatheory is explicit
that completeness here is relative to the normal-form grammar and the declared
observation relation, and does not establish that the relation contains every
physical timing, microarchitectural, operating-system, or network observation.

This is the sharpest available statement of what the architecture buys over a
flow-\emph{in}sensitive type system, which carries one fixed level per variable
and therefore rejects programs as simple as ``assign a secret, then overwrite it
with a constant, then publish.'' (A flow-sensitive system, with a per-program-point
environment, accepts that particular program --- it is their motivating example.
What no type system of either kind gets is the converse direction below.)
It costs no new proof --- only the discipline of stating a corollary that already
follows.

\subsubsection*{Where completeness actually leaks}

Every \Unknown{} the checker can return corresponds to a distinct formal reason
the biconditional's antecedent fails to be decidable:

\begin{center}
\begin{tabular}{@{}p{0.46\linewidth} p{0.44\linewidth}@{}}
\toprule
Formal cause & Reason-code family \\
\midrule
undecidability of reachability, so a bound is required
  & bound exhausted, resource limit \\
fragment membership: no semantics outside the whitelist
  & unsupported opcode, type, intrinsic, residual vector \\
missing \emph{relational} contract for a callee
  & missing contract, indirect call, recursion \\
solver incompleteness
  & solver timeout \\
the semantics is a \emph{set} of behaviours
  & freeze may choose, ambiguous NaN payload, uninitialized load \\
memory abstraction cannot decide aliasing
  & alias binding mismatch, layout-dependent pointer comparison \\
channel absent from the observation relation
  & unsupported address observation profile \\
\bottomrule
\end{tabular}
\end{center}

So the reason-code enumeration is not an error taxonomy. It is the exception list
of \cref{thm:complete}, which has a design consequence: the list must be
\emph{closed}, and each member owes an obligation of the form \emph{if this cause
does not apply, the checker decides}.

\begin{remark}[The half that is missing]
\label{rem:falsealarm}
The enumeration covers \emph{refusals}. It does not cover \emph{conservative
rejections} --- cases where the analysis reports a problem on a program that
satisfies the property. At least three such sources exist and none has a code:
over-approximated may-alias, an observation relation strictly stronger than
output-only noninterference, and the requirement of world-level constant
structure even at hosts no coalition can see. Because they carry no code, a false
alarm from any of them is indistinguishable in the report from a genuine finding.
If the completeness bookkeeping is to be as disciplined as the soundness
bookkeeping, these need first-class representation --- the same argument that
produced the precision controls of \cref{sec:pairs}, arriving from the theory
side rather than the testing side.
\end{remark}

A second, smaller theorem is worth mechanizing because it is where the
architecture could quietly go wrong:

\begin{theorem}[The diagnostic layer only narrows]
\label{thm:diag}
No diagnostic disposition implies $\Sec_\coal$, and no diagnostic disposition
removes a product constraint. In particular, a relational-required disposition
implies nothing whatever.
\end{theorem}

\cref{thm:diag} is small, and it is exactly the invariant that
\cref{rem:danger} says implementations break under pressure to reduce false
positives.

\subsubsection*{A milestone ladder}

\begin{enumerate}[leftmargin=1.8em]
\item \textbf{Mechanize the fragment's operational semantics}, with the
      observation projection as an explicit output rather than an afterthought.
      Deliverable: a step relation and a trace projection. The projection being
      explicit is what makes the property statable at all.
\item \textbf{State \eqref{eq:ni} and the product formally}, and prove
      \cref{thm:core}. This is the single highest-value proof in the programme,
      and it is independent of any implementation.
\item \textbf{Prove \cref{thm:diag}} for the actual disposition domain.
\item \textbf{Extract or validate the product encoder.} Either extract a verified
      encoder from the proof, or keep the fast one and check its output against
      the mechanized semantics per run. Prefer the latter.
\item \textbf{Do not mechanize the solver.} Consume unsatisfiability certificates
      from it and check those. Verifying an SMT solver is a research programme in
      its own right and buys little here.
\item \textbf{Validate the normalizer per run}, translation-validation style:
      show the normalized module is observationally equivalent to its input on
      this input, rather than proving the normalizer once. It will change
      constantly; the proof would not survive.
\item \textbf{Leave the backend to bounded validation.} Comparing the emitted
      machine code's observable trace structure against the IR-level one, on the
      actual artifact, is the only thing that addresses the
      trace-preservation obligation of \cref{step:deployment}. A verified backend
      would be better and is a far larger undertaking.
\end{enumerate}

\subsubsection*{What a proof can never give you}

Every \Lv{4} obligation is outside the reach of any proof about the IR: real
helper latency, whether the backend preserved the trace, hardware isolation
between tenants, speculative execution, sanitizer sufficiency, certificate
soundness. A mechanized theorem is always \emph{relative to a declared machine
model}, and the gap between that model and the silicon is exactly what the
\Conditional{} outcome exists to record.

This is worth stating loudly because it is the most common way formal claims
mislead. ``Proved noninterferent'' with an unstated machine model is not a
stronger claim than ``\Conditional{}, pending these four named facts'' --- it is
the same claim with the assumptions hidden. The four-outcome contract is, in that
light, less a limitation of this system than an unusually honest interface.

\subsection{Where this sits in the literature}
\label{sec:positioning}

Worth knowing precisely, because two things here are already owned and two are
not.

\paragraph{The product construction is classical.} Self-composition as a route to
noninterference is Barthe, D'Argenio and Rezk (CSFW 2004, extended in
\emph{MSCS} 2011); the framing of secure information flow as a safety problem is
Terauchi and Aiken (SAS 2005); the general theory of properties over sets of
traces is Clarkson and Schneider's hyperproperties (CSF 2008, \emph{JCS} 2010);
and relational Hoare logic is Benton (POPL 2004). For the type-system side of the
two-layer split, the canonical soundness result is Volpano, Irvine and Smith
(\emph{JCS} 4(2--3):167--187, 1996), surveyed by Sabelfeld and Myers
(\emph{IEEE JSAC} 2003).

\paragraph{Constant-time verification at LLVM level is occupied.} \emph{ct-verif}
(Almeida, Barbosa, Barthe, Dupressoir, Emmi; USENIX Security 2016) is the direct
ancestor: product construction, LLVM, annotations for public inputs. Two facts
about it matter. Its toolchain is effectively abandoned --- last commit 2017 ---
so it cannot be run as a baseline. But its benchmark corpus survives and is still
the field's de-facto suite: twelve directories covering \texttt{curve25519-donna},
\texttt{openssl}, \texttt{mee-cbc}, \texttt{polarssl}, \texttt{sodium},
\texttt{s2n} and others. \textbf{Reuse the corpus as fixtures; do not expect to
reuse the tool.} The live successor is \emph{CT-LLVM} (Zhang and Barthe, 2025),
which is analysis-based --- def-use plus alias analysis --- rather than
product-based, so it makes a different soundness/precision trade.

\paragraph{Compiler preservation is occupied.} The mechanized analogue of the
deployment obligation in \cref{step:deployment} is CompCert-CT, ``Formal
verification of a constant-time preserving C compiler'' (Barthe, Blazy,
Grégoire, Hutin, Laporte, Pichardie, Trieu; POPL 2020). The optimization-level
statement is Rosemann, Hack and Garg, ``Non-interference Preserving Optimising
Compilation'' (OOPSLA 2025), which develops hyperproperty simulations --- close
kin to the paired-refinement notion this system uses. The general framework for
what such a theorem should even say is Abate et al., ``Journey Beyond Full
Abstraction'' (CSF 2019).

\paragraph{What is not occupied.} Two gaps, and the second is larger.

First, a \emph{mechanized, product-construction} noninterference checker at LLVM
IR with a preservation story. The pieces exist separately --- Vellvm for the
semantics, the interaction-tree noninterference work for the property, CompCert-CT
for preservation --- and nobody has assembled them.

Second, and more striking: \textbf{MLIR has no information-flow story at all.}
There is no security, information-flow, taint, secrecy, or constant-time dialect
upstream. The nearest neighbour by name, HEIR's \texttt{secret} dialect, is a
lifting mechanism for homomorphic encryption; its own documentation states that
\texttt{secret.generic} does not itself do anything, and it carries no
noninterference property, no declassification, and no theorem. Mechanized MLIR
semantics exist, and information flow has not been built on them.

That is the opening, and it argues for a specific division of labour: prove the
core at LLVM-IR level where the semantics already exists, and treat MLIR as the
place where the \emph{policy surface} and the tensor-level channels live.

\subsection{If you only do three things}

\begin{enumerate}[leftmargin=1.8em]
\item Build the \textbf{exhaustive small-domain oracle} and fuzz against it. It
      finds real soundness bugs now, it produces regression fixtures as output,
      and it requires no proof assistant.
\item Make the checker \textbf{emit a certificate} and write the small
      independent validator. This shrinks what you have to trust, and it survives
      every future rewrite of the analysis.
\item Mechanize \textbf{\cref{thm:core} only}. One theorem, no implementation
      dependencies, and it is the statement everything else is a strategy for.
\end{enumerate}

Rungs 1--3 are weeks to months. A mechanized \cref{thm:core} over a small
fragment is months. A verified end-to-end pipeline is years, and would still
terminate in the same \Lv{4} obligations you started with.

\section{A checker that survives lowering}
\label{part:lowering}

Everything up to here analyses \emph{one} artifact. This part asks what happens
when that artifact is compiled further, and derives the architecture the
measurements force. The short version: single-point analysis is not salvageable,
and the fix is to stop pushing policy downward.

\subsection{What the measurements found}

Twenty semantics-preserving transformations were reproduced against LLVM 17.0.6.
They sort into three gaps, and it is worth naming them separately because they
fail in different directions.

\paragraph{The level gap.} The transformation happens \emph{below} LLVM IR, so an
IR-level verdict describes a program that does not run. The sharpest instance:
\code{x86-cmov-converter} turns a conditional move back into a branch. Composed
with mid-end speculation, which turns a branch into a select, the pipeline
produces the worst possible sequence --- \textbf{leaky source, clean \code{-O2}
IR, leaky binary}. The same IR yields opposite verdicts on x86-64,
aarch64-generic, aarch64-neoverse-n2, and riscv64. Related: \code{StackColoring}
merges two lifetime-bracketed allocations onto one frame offset, with
runtime-confirmed residue.

\paragraph{The annotation gap.} The policy does not survive. \code{combineMetadata}
discards custom metadata when merging; every signature-changing pass drops
parameter attributes; \code{mlir-translate} drops MLIR attributes entirely, with
no diagnostic. Combined with the convention that an unlabelled value is public,
each of these is a false negative rather than a refusal.

\paragraph{The counting gap.} Occurrence cardinality moves in \emph{both}
directions under plain \code{-O2} --- sinking merges two sites into one, jump
threading splits one into two. So no exact-count policy is an artifact
invariant.

\begin{remark}[The finding underneath all twenty]
\label{rem:notnormative}
The artifact the checker reads is not the artifact that runs, and nothing in the
system says which one is normative. Every other problem in this part is a
consequence.
\end{remark}

\subsection{Why single-point analysis cannot be repaired}

There are only two places to put a single analysis, and the measurements close
both.

\textbf{Analyse early}, where the policy is meaningful and the structure is
intact. Then the level gap applies in full: the binary can leak through a channel
that does not exist at that level, and the laundering sequence above produces a
clean verdict on a leaky program.

\textbf{Analyse late}, close to the machine, where the observations are real.
Then the annotation gap applies: the policy is gone. Worse, it is gone
\emph{silently}, so the analysis does not refuse --- it sees unlabelled code and,
under the usual convention, calls it public.

The obvious repair is to carry the policy downward so that late analysis has
something to read. \cref{sec:carriers} shows every known carrier has a silent-loss
path. That is what makes this an architectural problem rather than an engineering
one.

\subsection{Every carrier leaks}
\label{sec:carriers}

\begin{center}
\begin{tabular}{@{}p{0.30\linewidth} p{0.60\linewidth}@{}}
\toprule
Carrier & Measured failure \\
\midrule
MLIR \code{sps.*} attributes
  & \code{mlir-translate --mlir-to-llvmir} drops all of them; zero survive, no
    diagnostic \\
LLVM metadata
  & \code{combineMetadata} discards unknown kinds when merging two sites \\
Parameter attributes
  & dropped by every pass that rewrites a signature \\
Opaque marker call
  & must have no visible definition to survive, but then it does not link;
    once the definition is visible it is inlined away \\
Debug info
  & explicitly non-semantic; passes may drop or coarsen it \\
\bottomrule
\end{tabular}
\end{center}

Note what these have in common. They are all attempts to make a \emph{policy
fact} travel alongside the code, through a pipeline whose passes have no
obligation to preserve anything except semantics. The pipeline is not
misbehaving; it was never told these facts mattered.

\subsection{The inversion}
\label{sec:inversion}

If policy cannot travel down, send observations up.

\begin{center}
\begin{tikzpicture}[
  lvl/.style={draw, rounded corners=2pt, minimum height=7mm, minimum width=26mm,
              font=\small, inner sep=3pt},
  lbl/.style={font=\scriptsize},
  ar/.style={-{Latex[length=1.8mm]}}
]
  \node[lvl, fill=blue!8] (src) at (0,2.1)  {source / MLIR};
  \node[lvl]              (ir)  at (0,1.05) {LLVM IR};
  \node[lvl]              (mir) at (0,0)    {machine IR};
  \node[lvl]              (bin) at (0,-1.05){binary};

  \draw[ar] (src) -- (ir);   \draw[ar] (ir) -- (mir);   \draw[ar] (mir) -- (bin);

  \node[lbl, anchor=west, text width=4.3cm] at (1.9,2.1)
    {\textbf{policy lives here} --- authored, meaningful, stable};
  \node[lbl, anchor=west, text width=4.3cm] at (1.9,-0.5)
    {\textbf{observations live here} --- branches, addresses, frame slots, event order and count};

  \draw[ar, dashed] (3.4,-0.15) .. controls (5.2,0.9) .. (3.4,1.95);
  \node[lbl, anchor=west] at (4.6,0.95) {provenance};
\end{tikzpicture}
\captionof{figure}{Policy stays at the level where it is authored. Observations
are extracted at the level where they are real. A provenance map relates the
second to the first, and the relational check runs once, at the top, over
observations sourced from the bottom.}
\label{fig:inversion}
\end{center}

The three gaps close for structural reasons rather than by better engineering:

\begin{itemize}[nosep,leftmargin=1.4em]
\item \textbf{Level gap} --- observations are gathered where they are real, so a
      branch introduced by the backend is simply present in the observation set.
\item \textbf{Annotation gap} --- labels never travel, so nothing can drop them.
\item \textbf{Counting gap} --- cardinality is \emph{measured} at the bottom
      rather than asserted at the top, so a merged or duplicated site is
      observed, not assumed away.
\end{itemize}

\subsection{The quantifier error, and universal discharge}
\label{sec:quantifier}

The architecture as first stated has a defect worth dwelling on, because it is a
\emph{logic} error rather than an engineering gap, and because the first draft of
this section contained it.

That draft said: \emph{for each policy site, check the predicate over the
observations attributed to it}. That quantifies existentially over sites.
Soundness requires universal quantification over \emph{observations}.

The difference is not pedantic. A per-site predicate cannot distinguish
\textbf{``no observation here''} from \textbf{``compliant here''}. So every
mechanism that moves an observation off a named site --- outlining,
misattribution, a build without debug info --- leaves a site with nothing
attached, which the per-site rule reads as compliance. It \emph{fails open},
silently.

Reproduced with extractor, provenance map and artifact all held fixed: one
aarch64 build at \code{-Oz} yields \Verified{} under the per-site rule and a
refusal under the corrected one.

\paragraph{The correction.} \textbf{Every observation of a policy-relevant kind
in the artifact must be discharged by exactly one policy site.} Undischarged
observations are the failure condition; per-site predicates run only \emph{after}
that obligation is met.

The cost was measured, not estimated: across 42 fixtures $\times$ 2 targets
$\times$ 2 optimization levels, 167 of 168 builds admissible --- a refusal
surface near 1\%. It is affordable because it quantifies over policy-relevant
observation \emph{kinds}, not over instructions.

\begin{remark}[The same mistake, one level down]
The original problem was that policy could not survive travelling downward. The
first fix moved policy out of the way and rested the decision on the \emph{site
namespace} instead --- but a site namespace is another thing the compiler is free
to rewrite, so that repeats the original error against a different object.
Provenance is demoted accordingly: it may localize a finding, never decide one.
\end{remark}

\subsection{What each stage owes}

\begin{center}
\begin{tabular}{@{}l p{0.62\linewidth}@{}}
\toprule
Stage & Obligation \\
\midrule
Policy (top)
  & declare secrets, observer, releases, placement. Never travels down. It may be
    \emph{composed} across translation units: that is composition at the top and
    crosses no level \\
Extraction (each level)
  & emit the observation set for that level. This is what earns the
    architecture: observations existing at no higher level, such as an outlined
    function's secret-indexed loads, are seen only here \\
Provenance
  & a \textbf{localization hint}, not a decision mechanism. Three-valued:
    present, absent, or present-but-unattested. The third state is load-bearing
    --- confidently wrong provenance must degrade a diagnostic, never ship a
    leaky artifact marked verified \\
Decision (top)
  & \textbf{universal discharge}: every policy-relevant observation accounted for
    by exactly one site. Per-site predicates run afterwards \\
Differential
  & the actual soundness argument, not a supplement. Runs over \emph{whole
    images}, never over policy-named sites, under one composed program-level
    policy \\
Refusal
  & absent provenance, unattested provenance, absent callee model, or any
    undischarged observation \\
Identity
  & binds \emph{both} ends: input closure, link unit, target triple, subtarget
    features, codegen flags, compiler revision \\
\bottomrule
\end{tabular}
\end{center}

\begin{remark}[When this is worth the complexity, and when it is not]
\label{rem:simpler}
For a \emph{single pinned target tuple} a simpler architecture is equally sound
--- analyse the final artifact only, and re-supply policy at that level through a
side file keyed on symbol and offset rather than carried in the IR. At one site,
stack-slot merging, it is strictly \emph{more} sound.

The staged architecture earns its complexity only when more than one target tuple
ships. That is a checkable condition, and it should be checked before paying the
cost rather than assumed.
\end{remark}

\begin{remark}[The role of markers changes, and shrinks]
Under \cref{sec:inversion} a declassification marker no longer has to survive to
the binary. It only has to survive to the level where the \emph{policy} is
evaluated. That is a far weaker requirement than the one the attack findings
demolished, and it is why the inversion is worth the engineering: it converts a
problem with no known sound solution into one with a workable one.
\end{remark}

\subsection{The declassification marker, after the attacks}
\label{sec:marker}

Markers still have a job: recording that a High-to-Low transition was
\emph{authorized}. The obvious realization --- an opaque function the optimizer
cannot see through --- was attacked across five dimensions and does not survive.
The failure is structural, not a tuning problem.

\paragraph{Why the function-shaped marker fails.} A declaration must resolve at
link time. The moment its definition is co-visible --- LTO, ThinLTO, a header,
the same translation unit --- attribute inference deduces \code{returned} from
the identity body, the optimizer replaces uses of the marker's result with its
argument, and the release edge is severed \emph{while the marker sits there
decoratively}. A presence-based checker sees a marker and a clean path; the raw
secret reaches the sink. Verified end-to-end through the real \code{clang} driver
under full LTO.

That is the whole family: no definition means no link, and a definition means no
protection.

\paragraph{What works instead: tied-register inline asm.} There is no symbol, so
there is nothing to define, nothing to link, and nothing for LTO to import.

\begin{center}
\begin{minipage}{0.92\linewidth}
\small
\begin{verbatim}
%h = call T asm sideeffect "// sps.declassify id=$2",
                           "=r,0,i,~{memory}" (T %raw, i32 <ID>)
\end{verbatim}
\end{minipage}
\end{center}

Each operand earns its place, and each was measured:

\begin{center}
\begin{tabular}{@{}l p{0.68\linewidth}@{}}
\toprule
Operand & What it buys \\
\midrule
\code{sideeffect}
  & prevents deletion when the result is unused, and prevents machine-level CSE
    merging two markers \\
\code{=r} output
  & puts the marker \emph{on} the dataflow path. With no body, \code{returned}
    cannot be inferred --- the exact failure of the function-shaped form \\
\code{0} tie
  & ties input to output, so the machine cost is zero. Measured 9 instructions
    marked versus 9 unmarked, no frame, no spills, tail call preserved \\
\code{i} on the id
  & the id must be an immediate. A non-constant id is a \textbf{hard
    instruction-selection error}, so identity recoverability is fail-closed by
    construction rather than by convention \\
\code{\textasciitilde\{memory\}}
  & no IR-level effect, but blocks post-register-allocation scheduling from
    reordering memory operations across the marker. Measured cost zero \\
\bottomrule
\end{tabular}
\end{center}

Verified identical under \code{default<O1/O2/O3/Os/Oz>}, \code{lto<O2>},
\code{lto<O3>} and \code{thinlto<O2>}: marker present, id literal, sink consuming
the marker's result --- while the declassifier implementation is fully inlined
and deleted.

\paragraph{The freeze point this forces.} Post-LTO, post-optimization IR,
immediately before instruction selection. The asm comment is dropped by the
assembler --- measured one occurrence in the \code{.s}, zero in the \code{.o},
zero in the linked binary --- so anything later cannot see the marker at all.
Where a shipped-binary attestation is genuinely required, the comment body can be
replaced by a section note emitting a \code{(pc, release\_id)} record with zero
text instructions.

\paragraph{Two problems that do not go away.} Stated as residual risk rather than
buried:

\begin{itemize}[nosep,leftmargin=1.4em]
\item \textbf{\code{returned} is still expressible.} Inline asm does not
      structurally forbid it --- the IR verifier accepts \code{returned} on an
      asm argument, and the bypass reproduces. It must be an explicit checker
      rule, because there is no IR-level prohibition.
\item \textbf{A static marker count is not a release count.} Unrolling and
      inlining both multiply it. Occurrence cardinality has to be restated
      dynamically, or checked at an earlier freeze point --- which is the
      counting gap of \cref{rem:notnormative} reappearing, unsolved.
\end{itemize}

\begin{remark}[The general lesson]
The function-shaped marker failed because it asked the \emph{linker} and the
\emph{optimizer} to respect a policy fact neither was told about. The inline-asm
marker works because it asks for nothing: it is opaque by construction, not by
permission. Every durable mechanism in this part has that shape --- it survives
because there is no decision point at which something could decide otherwise.
\end{remark}

\subsection{The differential check}

Even with the inversion, run the analysis at each level independently and compare
verdicts. Disagreement between levels is itself the finding:

\begin{center}
\begin{tabular}{@{}l l p{0.44\linewidth}@{}}
\toprule
Upper & Lower & Meaning \\
\midrule
\Verified{} & \Unsafe{}   & \textbf{laundering} --- the lowering introduced the leak.
                            No fixture in the corpus covers this direction \\
\Unsafe{}   & \Verified{} & evidence destroyed; the leak was optimised out of view \\
\Verified{} & \Unknown{}  & acceptable: monotone degradation \\
\Unknown{}  & \Verified{} & suspicious --- what fact appeared? \\
\bottomrule
\end{tabular}
\end{center}

This is cheap. It needs no new theory, it reuses the existing corpus, and it is
the mechanism that would have caught the inlined-declassifier case automatically
rather than by inspection.

\subsection{What this costs, and what it does not buy}

The honest accounting.

\textbf{It costs backend work.} Extraction at machine-IR level is not an IR pass;
it means reading or instrumenting the target layer, per target. The measurements
already show four target tuples disagreeing on one input, so this is per-tuple
work.

\textbf{Provenance is best-effort.} Debug-info \code{inlinedAt} chains do relate
inlined instructions back to their original function --- verified --- but debug
info is explicitly non-semantic and may be dropped. So the provenance map is a
source of \Unknown{}, not a guarantee. That is acceptable precisely because
\Unknown{} is a first-class outcome; it would be unacceptable in a three-valued
design.

\textbf{It is not a proof.} This architecture makes the checker's answer be about
the artifact that runs. It does not establish that the answer is right --- that
is still \cref{thm:core} and the assurance ladder of
Part~\ref{part:correctness}.

\textbf{It does not close \Lv{4}.} Real latency, hardware isolation, and
speculation remain outside any observation set you can extract statically.

\subsection{One capture point, several extraction points}
\label{sec:capture}

Two results in this document appear to disagree about where the artifact is
frozen, and the disagreement is terminological. Resolving it fixes the most
important architectural parameter, so it is worth doing explicitly.

The marker analysis concludes the freeze point is \textbf{post-LTO, immediately
before instruction selection}, because the marker is not visible below that ---
measured, one occurrence in the \code{.s}, zero in the \code{.o}, zero in the
linked binary. The inversion concludes observations must come from \textbf{the
machine level}, because that is where \code{je} rather than \code{csel} is
decided. Both are right. They are answers to different questions that were being
given one name.

Split the concept:

\begin{center}
\begin{tabular}{@{}l p{0.62\linewidth}@{}}
\toprule
\textbf{Capture point} & where the \emph{annotated} artifact is snapshotted and
hashed. Exactly one. Must be where policy is still readable: post-LTO, before
instruction selection. Markers must survive to here and no further. \\[3pt]
\textbf{Extraction points} & where \emph{observations} are gathered. Several, at
and below the capture point, including machine IR. Carry no policy at all. \\[3pt]
\textbf{Artifact identity} & binds the capture point to every extraction point,
so the report describes one compilation rather than a family. \\
\bottomrule
\end{tabular}
\end{center}

The requirement on markers weakens accordingly, and this is the practical
payoff: a marker must survive to the capture point, not to the binary. That is a
far easier obligation, and it is why \cref{sec:marker}'s inline-asm form is
sufficient despite the assembler discarding it.

\begin{remark}[What the identity must bind]
Binding the top-level module alone is not enough --- the same module produced
\code{je} on x86-64 and \code{csel} on aarch64 from identical IR. So the identity
is a \emph{tuple}: module hash, ordered pipeline digest, target triple, target
CPU and features, and optimization level. A verdict that does not name that tuple
is not a claim about anything runnable. This is the same shape of requirement as
per-coalition rows, and it lands on the same record.
\end{remark}

\subsection{Open problems}

\paragraph{Irreducible: event multiplicity.} This one is not an open problem to
be solved later; it is a limit of the architecture, and naming it is the only
honest handling.

Provenance can be \emph{perfect} and the count still wrong. Measured: at
\code{-O2 -g}, two \code{stp} store-pair instructions carrying entirely correct
source attribution perform sixteen release events --- vectorization at four to
one, composed with store-pairing at two to one. Nothing is misattributed. The
provenance map is exactly right. The count is off by a factor of eight.

Worse, the defect is structural rather than incidental. The decision quantifies
over the \emph{union} of observation sets, and set union is idempotent, so the
formulation cannot count at all. A policy constraining release cardinality is
therefore not checkable in this architecture as stated. It needs either a
multiset, which changes the decision procedure, or a dynamic statement, which
changes what is being claimed.

Until one of those is chosen, a release-count policy must be reported
\Unknown{}, not approximated.

\medskip
Two further problems remain open, stated rather than hidden.

\textbf{Non-monotone cardinality.} Because occurrence counts move both ways, a
policy that constrains release count cannot be checked as an artifact invariant
at a single level. It has to be a statement about the observation set, which
means the extraction has to be trusted to be complete.

\textbf{Multi-target verdicts.} One artifact and four targets gave four different
answers. Either a verdict names its target tuple, or the artifact must be
verified for every tuple it will be built for --- and the report needs a row per
tuple, which is the same record-shape problem as coalitions.

\section{From design to implementation}
\label{part:implementation}

Every preceding part describes something that does not exist. This part is about
the one piece that does, what it measured, and why the measurement matters more
than the six parts before it.

\subsection{What actually exists}

It is worth being exact, because the distinction between a specification and a
system is the easiest one to lose while writing a document like this.

\begin{center}
\begin{tabular}{@{}l r p{0.44\linewidth}@{}}
\toprule
Artifact & Size & Kind \\
\midrule
MLIR fixtures            & 47   & test data \\
C provenance files       & 42   & test data \\
Actor examples           & 7    & test data \\
Integration tests        & 8    & test data \\
\midrule
\code{check\_harness.py} & 929  & metadata \emph{validation}; checks fixtures are
                                  well-formed. Not analysis \\
\code{refusal\_rate.py}  & 179  & a scanner over IR text. Not analysis \\
\code{sps-scan.cpp}      & 174  & \textbf{an actual analysis}: an MLIR
                                  \code{SparseDataFlowAnalysis} with a
                                  \textsf{Low}/\textsf{High} lattice \\
\bottomrule
\end{tabular}
\end{center}

So: roughly a hundred test files, and of thirteen hundred lines of code exactly
one hundred and seventy-four implement an information-flow analysis. Everything
else validates, measures, or documents.

\begin{remark}[Why the accounting matters]
It is easy to write ``the architecture refuses 63\% of functions'' when what was
measured is ``63\% of functions contain an operation with no source location.''
The first is a claim about a system; the second is a claim about clang's output
plus a rule from a design document. Only the second was measured. Keeping the
two apart is the difference between a result and a rationalisation.
\end{remark}

\subsection{The first real measurement}

\code{sps-scan} was run over every fixture whose policy survives parsing --- 17
of 47, the rest encoding their policy in SSA names that \code{mlir-opt}
discards. It emits findings with a stable reason and a source location:

\begin{center}
\begin{minipage}{0.9\linewidth}\small
\begin{verbatim}
[secret-dependent-branch] llvm.cond_br @ loc(...:67:5)
[secret-to-public-sink]   llvm.store   @ loc(...:69:5)
findings: 2
\end{verbatim}
\end{minipage}
\end{center}

The results sort into four groups, and all four are informative.

\paragraph{Caught, correctly.} \fx{predecessor\_choice\_blockarg.bad} yields two
findings, including the MT-CM6 case where a secret selects a merge block argument
with no secret on any SSA operand edge. \fx{prefix\_causal\_release.bad} yields
one. These are the countermodels a re-implementation would most plausibly fail,
and a 174-line pass catches them.

\paragraph{Silent, correctly.} Five fixtures yield zero findings: the
public-allocation control, the offset-disjoint control, the overwritten-slot
control, the repaired error oracle, and the constant-latency wolfSSL helper
profile. A pass that reported everything as unsafe --- the failure mode the
corpus was previously unable to detect at all --- fails here.

\paragraph{Imprecise, as predicted.} Two precision controls fire, one finding
each: the identical-successor branch and the value-cancelling \code{xor}.

This is \textbf{not} a defect. Those fixtures exist to document that a forward
\textsf{Low}/\textsf{High} lattice \emph{must} lose precision at exactly those
sites, and each carries \code{l1 disposition: relational-required} saying so. The
implementation is behaving the way \cref{step:diag} says it is required to. It is
the first time a prediction in this document was confirmed by running something.

\paragraph{Missed, for known reasons.} Three bad fixtures yield zero: the
audience mismatch, the error oracle, and the affected wolfSSL helper profile. They need audience reasoning, release-policy conformance and a
helper-latency contract respectively --- none of which a taint pass implements.
These are the shape of the remaining work, and they are visible only because
something ran.

\subsection{The argument for building before designing further}

This document has been wrong, and corrected by measurement, six times:

\begin{center}
\begin{tabular}{@{}p{0.43\linewidth} p{0.45\linewidth}@{}}
\toprule
Claim & Corrected by \\
\midrule
whole-run release equality as the
    product's initial constraint
  & the spec forbids it; the encoding
    is unsound in the permissive
    direction \\
the declassification marker as an
    opaque function
  & attribute inference deducing
    \code{returned} once the definition
    is co-visible \\
per-site quantification in the
    decision rule
  & a build reporting \Verified{} under
    it and refusing under universal
    discharge \\
a register-only reduction of the laundering case
  & it did not reproduce; the conversion needs a memory operand \\
an arithmetic-mask repair
  & folded straight back into a select, byte-identical to the leaky twin \\
a refusal surface of about 1\%
  & 4.8\% per observation, and 63\% per function on real code \\
\bottomrule
\end{tabular}
\end{center}

Six for six. Not one of these was caught by review, by internal consistency, or
by re-reading the specification; every one was caught by running a tool and
looking at the output. The corrected versions are better, but the pattern is the
point: \textbf{in this problem domain, a design claim that has not been executed
is not evidence.}

The remaining untested claims are all design claims. So the highest-value next
action is not more design --- it is to make the smallest real thing falsifiable.

\subsection{The smallest useful next step}

\code{sps-scan} is not referenced anywhere in the harness. Wiring it in converts
the whole corpus from a specification into a test of an implementation:

\begin{enumerate}[nosep,leftmargin=1.6em]
\item Add it to the diagnostic \code{RUN} of the satisfiable bring-up gates. That
      turns ``the analysis must emit this reason here'' from an aspiration into a
      pass/fail.
\item Add it to the four precision controls under their declared
      \RelReq{} disposition. This is the first false-positive regression net the
      corpus has ever had --- and the run above shows exactly which two controls
      it will legitimately trip.
\item Re-derive the refusal number from the analysis rather than from a text
      scanner, with a real policy scoping the observation set.
\end{enumerate}

None of this needs the lowering architecture of Part~\ref{part:lowering}, the
\code{.ll}-versus-MLIR decision, or the thirty unmigrated fixtures. Those gate
later layers. The diagnostic layer is buildable now, against a corpus that is
already known to exercise it.

\begin{remark}[What is durable]
The architecture in this document has been rewritten twice in a day and will
likely change again. What has not changed, and will survive any rewrite, is the
set of reproduced compiler behaviours, the fixtures that encode them, and the
measurement scripts. Those are the assets. A design that cannot yet be executed
should be weighted accordingly --- including this one.
\end{remark}

\appendix
\section{Appendix: code map, report shape, and references}
\label{part:appendix}

\subsection{Where things live}

All paths below are relative to \fx{prototypes/compiler\_harness/} in the
repository root.

\begin{center}
\begin{tabular}{@{}p{0.30\linewidth} p{0.60\linewidth}@{}}
\toprule
Path & Contents \\
\midrule
\fx{fixtures/}
  & 54 case directories, each pairing one MLIR shape with its snapshot; family
    directories also own C evidence and optional case-local candidate bundles \\
\fx{c/}
  & shared build/check support and the executable equivalence driver \\
\fx{c/check\_harness.py}
  & structural validation for snapshots, evidence paths, and case-local
    candidates \\
\fx{integration/}
  & metadata, witness-driver, import, and code-generation tests \\
\fx{integration/candidate-bundles/}
  & reproducibility and negative tests for the nine nonclaimable candidates \\
\fx{examples/actors/}
  & actor and ACL design examples; outside the enforced corpus \\
\fx{FIXTURE\_REVIEW\_GUIDE.md}
  & review order, fixture intent, corpus counts, and manual sign-off criteria \\
\bottomrule
\end{tabular}
\end{center}

Two orientation notes for a reader opening the corpus for the first time. A
case's sibling \fx{snapshot.yaml} is the authoritative preflight boundary and
expectation; its \fx{c\_evidence} paths record provenance, not a claim that the
MLIR was compiled from that C. The fixture inventory is \emph{closed}: every
MLIR case has exactly one sibling snapshot and every snapshot has exactly one
sibling MLIR file, while every family C source is cited by at least one case.

\subsection{Running it}

\begin{center}
\begin{tabular}{@{}p{0.42\linewidth} p{0.5\linewidth}@{}}
\toprule
Command & Effect \\
\midrule
\fx{make check}            & all IR and integration tests \\
\fx{make check-shape}      & all 54 fixture MLIR/snapshot shape tests \\
\fx{make check-mlir}       & compatibility alias for \fx{check-shape} \\
\fx{make check-candidates} & the nine candidate-bundle tests \\
\fx{make check-artifacts}  & compatibility alias for \fx{check-candidates} \\
\fx{make check-integration}& metadata, witness, import, code generation \\
\fx{make bootstrap-lit}    & create the pinned local test runner \\
\bottomrule
\end{tabular}
\end{center}

A clean default suite is the correct outcome: executable preflight checks pass,
while future SPS semantic tests remain unsupported until a Rev-4 verifier and
canonical materialized bundle are supplied. Neither a green preflight test nor
an unsupported semantic arm computes \fx{ModelStatus}.

\subsection{The shape of a report}

The walkthrough implies a record with more structure than a single verdict token.
Collecting what each step required:

\begin{center}
\begin{tabular}{@{}l p{0.60\linewidth}@{}}
\toprule
Field & Constraint that forces it \\
\midrule
model status
  & \Proved{} carries no reason; a counterexample carries a witness; a refusal
    carries an ordered blocker set (\cref{step:verdict}) \\[3pt]
deployment status
  & an independent axis: a model counterexample cannot be repaired by deployment
    evidence (\cref{step:deployment}) \\[3pt]
policy review status
  & a third axis that by construction never changes either of the other two \\[3pt]
diagnostic disposition
  & per site, from the five-valued domain of \cref{step:diag}, plus a health
    flag; never a premise for safety \\[3pt]
row key
  & one row per entry and coalition, over the derived closure
    (\cref{step:coalitions}) \\[3pt]
witness
  & replayable, and canonical enough that the discovery route is not
    observable in it (\cref{step:verdict}) \\
\bottomrule
\end{tabular}
\end{center}

The three-status split is the part most likely to be collapsed under time
pressure, and \Conditional{} is precisely the projection of ``model proved,
deployment open.'' It is a derived value rather than a primitive alongside the
other three statuses --- which is why
an obligation naming a target fact and one naming a backend fact should not share
a field.

\subsection{References}
\label{sec:refs}

Verified against Crossref, dblp, and author- or proceedings-hosted copies.
Publisher pages for USENIX, IEEE, ACM, and Springer were not directly
retrievable, so those entries rest on the registration authority rather than the
publisher's own page.

\begingroup
\small
\begin{description}[leftmargin=1.2em,style=unboxed,itemsep=2pt]

\item[Self-composition and 2-safety.]
G.~Barthe, P.~R. D'Argenio, T.~Rezk. \emph{Secure Information Flow by
Self-Composition}. CSFW 2004, 100--114; extended in \emph{Math.\ Struct.\ in
Comp.\ Sci.}\ 21(6):1207--1252, 2011.\quad
T.~Terauchi, A.~Aiken. \emph{Secure Information Flow as a Safety Problem}. SAS
2005, 352--367.\quad
M.~R. Clarkson, F.~B. Schneider. \emph{Hyperproperties}. CSF 2008, 51--65;
\emph{J.\ Computer Security} 18(6):1157--1210, 2010.\quad
N.~Benton. \emph{Simple Relational Correctness Proofs for Static Analyses and
Program Transformations}. POPL 2004, 14--25.

\item[Information-flow type systems.]
D.~M. Volpano, C.~E. Irvine, G.~Smith. \emph{A Sound Type System for Secure Flow
Analysis}. \emph{J.\ Computer Security} 4(2--3):167--187, 1996.\quad
A.~Sabelfeld, A.~C. Myers. \emph{Language-Based Information-Flow Security}.
\emph{IEEE JSAC} 21(1):5--19, 2003.

\item[Constant-time verification.]
J.~B. Almeida, M.~Barbosa, G.~Barthe, F.~Dupressoir, M.~Emmi. \emph{Verifying
Constant-Time Implementations}. USENIX Security 2016, 53--70. \emph{(ct-verif;
toolchain unmaintained since 2017, benchmark corpus still current.)}\quad
Z.~Zhang, G.~Barthe. \emph{CT-LLVM: Automatic Large-Scale Constant-Time
Analysis}. IACR ePrint 2025/338.\quad
L.-A. Daniel, S.~Bardin, T.~Rezk. \emph{Binsec/Rel}. IEEE S\&P 2020, 1021--1038;
extended in \emph{ACM TOPS} 26(2), 2023.\quad
L.~Cai, F.~Song, T.~Chen. \emph{Towards Efficient Verification of Constant-Time
Cryptographic Implementations}. FSE 2024.

\item[Compiler preservation of hyperproperties.]
G.~Barthe, S.~Blazy, B.~Gr\'egoire, R.~Hutin, V.~Laporte, D.~Pichardie,
A.~Trieu. \emph{Formal Verification of a Constant-Time Preserving C Compiler}.
\emph{PACMPL} 4(POPL), Art.~7, 2020.\quad
J.~Rosemann, S.~Hack, D.~Garg. \emph{Non-interference Preserving Optimising
Compilation}. \emph{PACMPL} 9(OOPSLA2):1429--1454, 2025.\quad
C.~Abate, R.~Blanco, D.~Garg, C.~Hri\c{t}cu, M.~Patrignani, J.~Thibault.
\emph{Journey Beyond Full Abstraction}. CSF 2019, 256--271.\quad
M.~Patrignani, A.~Ahmed, D.~Clarke. \emph{Formal Approaches to Secure
Compilation}. \emph{ACM Computing Surveys} 51(6), 2019.

\item[Mechanized LLVM and MLIR.]
J.~Zhao, S.~Nagarakatte, M.~M.~K. Martin, S.~Zdancewic. \emph{Formalizing the
LLVM Intermediate Representation for Verified Program Transformations}. POPL
2012, 427--440.\quad
Y.~Zakowski, C.~Beck, I.~Yoon, I.~Zaichuk, V.~Zaliva, S.~Zdancewic.
\emph{Modular, Compositional, and Executable Formal Semantics for LLVM IR}.
\emph{PACMPL} 5(ICFP), 2021. \emph{(current Vellvm.)}\quad
L.~Silver, P.~He, E.~Cecchetti, A.~K. Hirsch, S.~Zdancewic. \emph{Semantics for
Noninterference with Interaction Trees}. ECOOP 2023, 29:1--29:29.\quad
N.~P. Lopes, J.~Lee, C.-K. Hur, Z.~Liu, J.~Regehr. \emph{Alive2: Bounded
Translation Validation for LLVM}. PLDI 2021, 65--79.\quad
M.~Fehr, Y.~Fan, H.~Pompougnac, J.~Regehr, T.~Grosser. \emph{First-Class
Verification Dialects for MLIR}. \emph{PACMPL} 9(PLDI):1466--1490, 2025.\quad
S.~Bhat, A.~Keizer, C.~Hughes, A.~Goens, T.~Grosser. \emph{Verifying Peephole
Rewriting in SSA Compiler IRs}. ITP 2024, 9:1--9:20. \emph{(lean-mlir.)}

\item[Constant-time by construction.]
J.~B. Almeida et al. \emph{Jasmin: High-Assurance and High-Speed Cryptography}.
CCS 2017, 1807--1823; \emph{The Last Mile}, IEEE S\&P 2020, 965--982.\quad
S.~Cauligi et al. \emph{FaCT: A DSL for Timing-Sensitive Computation}. PLDI
2019, 174--189.

\end{description}
\endgroup

\subsection{Reading the arc backwards}

One last framing, for a reader who wants the shortest possible summary of
Part~\ref{part:walkthrough}.

The checker never proves a program safe. It reduces ``is this program safe'' to
``is this product unreachable,'' having first established that the question is
well-posed: the policy is total, the artifact is in the fragment and is the one
that shipped, and the admitted domain is non-empty. Every one of those
preconditions is a place where a plausible shortcut yields a confident wrong
answer, and the corpus exists because each shortcut is more tempting than it
looks.


\end{document}
